\documentclass[a4paper,12pt,openany,twoside, abstracton]{book}%classe report di KOMA-Script
\usepackage[utf8]{inputenc} %accenti italiani
\usepackage[parts,pdfspacing,eulerchapternumbers,eulermath,dottedtoc]{classicthesis} %parts,pdfspacing,eulerchapternumbers,eulermath,dottedtoc
\usepackage{arsclassica}
\usepackage[english]{babel} % serve a far si che i capitoli vengano scritti in italiano
\usepackage[utf8]{inputenc} % serve per scrivere direttamente i caratteri accentati
\usepackage[T1]{fontenc} 
\usepackage{microtype} 
\usepackage{graphicx} % serve per le
\usepackage{amssymb}
\usepackage{mathrsfs}
\usepackage{framed}
\usepackage{bm}
\usepackage{fancybox}
\usepackage{textcomp} %per inserire °
\usepackage{colortbl} %serve per colorare le tabelle
\usepackage{color}
\usepackage{hyperref} %[colorlinks=false]
\usepackage{amsfonts}
\usepackage{listings} %per l'inserimento di codice
\usepackage{amsmath} %testo dentro equazioni
\usepackage{mathrsfs}
\usepackage{graphicx}
\usepackage{subfig} %matrice di immagini
\usepackage[numbers]{natbib} %bibliografia
\usepackage[tight,italian]{minitoc} % mini indice
\usepackage{rotating}
\usepackage{caption}
\usepackage{enumerate}
\usepackage{graphicx}
\usepackage{floatflt}
%\usepackage{biblatex}
%\usepackage{subfigure}
\captionsetup{tableposition=top,figureposition=bottom,font=small}

\usepackage{tikz}
%\usepackage{fullpage}

\usepackage[a4paper,top=2.5cm,bottom=2.5cm,left=3cm,right=3cm,bindingoffset=1mm]{geometry}
\usepackage{multirow}
\usepackage{stackengine}
\usepackage{lmodern}
\usepackage{setspace}

\usetikzlibrary{arrows,shapes,positioning,shadows,trees}

\tikzset{
	basic/.style  = {draw, text width=3cm, drop shadow, font=\sffamily, rectangle},
	root/.style   = {basic, rounded corners=2pt, thin, align=center,
		fill=MidnightBlue!80},
	level 2/.style = {basic, rounded corners=2pt, thin,align=center, fill=ForestGreen!30,
		text width=8em,drop shadow,},
	level 3/.style = {basic, thin, align=center, fill=YellowGreen!15, text width=6.5em, rectangle}
}

%$\frac{\tau\pm\epsilon \cdot V}{\%12}$ esempio di formula
%chapter
%section
\graphicspath{{images/}}

\makeatletter
\newcommand*{\lolnoheading}{\@starttoc{lol}}
\makeatother

\begin{document}
\fontfamily{ptm} \selectfont
%Frontespizio



\begin{titlepage}
\begin{center}
% Upper part of the page
 


{{\Large{\textsc{Politecnico di Torino}}}} 
\rule[0.1cm]{\textwidth}{0.1mm}
\rule[0.5cm]{\textwidth}{0.6mm}

{\small{\bf CORSO DI LAUREA MAGISTRALE IN INGEGNERIA ELETTRONICA}}
\end{center}
\vspace{5mm}
\begin{center}
	\includegraphics[width=0.3\textwidth]{./Frontespizio/logo_polito}\\[1cm] 
\end{center}


\vspace{10mm}

\begin{center}
	 {\large{Tesi di Laurea Magistrale}}\\
	 \vspace{5mm}
{\LARGE{\bf Google Glass Data Visualization and}}\\
\vspace{1mm}
	{\LARGE{\bf Monitoring for Organs-on-a-Chip and}}\\
	\vspace{1mm}
	{\LARGE{\bf Biomedical Applications}} \\ 

\vspace{15mm}
\end{center}
{\large
	Relatore\\
		{
			\bf 
			 Prof. Danilo Demarchi\\ 
			
		}
}
\begin{flushright}
	{\large
		Candidato\\
		{
			\bf 
			Fabio Busignani s197883\\ 
			
		}
	}

\end{flushright}

\vfill
\begin{center}
{\large Marzo 2015}
\end{center}


\end{titlepage}
\newpage
\pagenumbering{Roman}


\selectlanguage{english}
\linespread{1.5}
%\begin{abstract}
\chapter*{Abstract}
\addcontentsline{toc}{chapter}{Abstract}

\begin{figure}[h]
	\subfloat{%First sub-figure\label{subfig-1:dummy}]{%
		\includegraphics[width=.25\textwidth]{./firstpage/MIT2}
	}
	\hfill
	\subfloat{%}First sub-figure\label{subfig-2:dummy}]{%
		\includegraphics[width=.25\textwidth]{./firstpage/Harvard}
	}
	\subfloat{%}First sub-figure\label{subfig-2:dummy}]{%
		\includegraphics[width=.25\textwidth]{./firstpage/BWH}
	}
	\subfloat{%}First sub-figure\label{subfig-2:dummy}]{%
		\includegraphics[width=.25\textwidth]{./firstpage/HST}
	}
\end{figure}


The present scripture represent the master thesis of Fabio Busignani, and it has been carried out at \href{http://www.tissueeng.net/lab/}{\textit{Khademhosseini Lab}} (Cambdridge, MA, USA), \href{https://hst.mit.edu/}{Harvard-MIT Health Science and Technology}, \href{http://www.brighamandwomens.org/}{Brigham and Women's Hospital} under the supervision of professor Ali Khademhosseini and Ph.D. Yu Shrike Zhang.\\
The design which is going to be described, has been inserted inside the context of a five years project (\textit{\textbf{XCEL}} grant), sponsored by the U.S. Defense Threat Reduction Agency (\textit{DTRA}).\\

The aim of \textit{XCEL} is to develop a \textit{Body-On-A-Chip} microfluidic platform that is able to simulate multi-tissue interactions under physiological fluid flow conditions.\\

This master thesis will focus on designing a custom user interface on \textit{Google Glass} for simultaneous recording of biosensing data such as temperature, pH, and microscopy images/videos as well as remote control of microfluidic valves and devices.  The project involves all the hierarchical layers, starting from the physical one with the circuit in charge to acquire data from bio-sensors and drive the valves, up to the glasswear\footnote{Google Glass Application}.\\

In the Introduction chapter, the main keys of the project are presented in detail as well as the final result from a user point of view.\\
After that, a detailed description of each abstraction level which goes to build the entire systems is shown: starting from the bottom (Hardware) reaching the top (Google Glass Application) passing through the Firmware, that runs in an embedded \textit{Linux} platform, and the Software, present on the \textit{PC} and the \textit{Google App Engine}.\\
The  Experiments and Conclusion chapter shows the obtained results with different experiments. The thesis ends discussing the limitations and improvements of the system.



\cleardoublepage
\tableofcontents
\cleardoublepage
\pagenumbering{arabic}

\chapter*{INTRODUCTION}\label{summary}
\addcontentsline{toc}{chapter}{{\footnotesize INTRODUCTION}}
\markboth{Fabio Busignani Master Thesis}{INTRODUCTION}
This master thesis has been carried out at Khademhosseini laboratory, Harvard-MIT Health Science and Technology (Brigham and Women's Hospital), in Cambridge, MA.\\
During my six-months of research I have joined \textit{\textbf{XCEL}} grant project, a five years project sponsored by the U.S. Defense Threat Reduction Agency (\textit{DTRA}).

The goal of this project is to develop a system, a microscale bioreactor containg four 3D fully-functional \textit{organs-on-a-chip}.

\begin{figure}[h]
	\centering
	\includegraphics[width = 0.8\textwidth]{Intro/body}
	\caption{\textit{XCEL} project (\textit{Body-on-a-chip})}
	\label{Fig:Body}
	
\end{figure}

The (Fig.\ref{Fig:Body}) shows, in a schematic representation, the design of the entire \textit{XCEL} project. On the breadboard four organs are connected from each others: liver, heart, ling, and vascular system (\textit{Vasc}). The medium used to connect them is using tubing circuit where the media flows. This tubing connections are driven by electrovalves.


My role in this project has been to create a custom user interface on Google Glass for simultaneous recording of biosensing data such as temperature, pH, and microscopy images/videos as well as remote control of the microfluidic valves previously introduced. In summary my aim was to design a Google Glass App for use in organs-on-a-chip platforms.

The \textit{organs-on-a-chip} platforms contain interconnected microfluidic modular components including the bioreactors for hosting biomimicry human organ models, downstream biochemical sensors to continually monitor the levels of biomarkers secreted by the organs, and physical sensors to monitor the physical microenvironment of the circulatory system. Due to their extensive similarity with human organs, these miniature human models are finding widespread applications where the prediction of in vivo responses of the human body is needed, including but not limited to drug screening, basic biomedical studies, and environmental safety assessment. Thus, the \textit{organs-on-a-chip} platforms seek to recapitulate human organ function at micro-scale by integrating microfluidic networks with three-dimensional organ models, which are expected to provide robust and accurate predictions of drug/toxin effects in human bodies. In fulfilling this aim, a set of physical/chemical parameters need to be monitored and stored in order to capture such effects of drug/toxin administered into the system.\\

By precisely designing the Google Glass App for this organs-on-a-chip platform it allows convenient observation and control of the organ models, biosensors, and the microfluidic circuitry, which has been difficult to achieve previously.

The system designed and described in this thesis is illustrated in (Fig.\ref{Fig:BlockDiagram}).
 
 \begin{figure}[h]
 	\centering
 	\includegraphics[width = \textwidth]{Intro/Block_Diagram.eps}
 	\caption{Block diagram of the system}
 	\label{Fig:BlockDiagram}
 	
 \end{figure}
 
 The (Fig.\ref{Fig:BlockDiagram}) shows the principal steps of data transmission from physical and video sensors to the Google Glass via an \textit{Embedded Linux System} performed using the \href{http://beagleboard.org/BLACK}{\textbf{Beaglebone Black}}.\\
 The Beaglebone Black runs processes that are in charged to:
 \begin{itemize}
 	\item acquire the sensors value and to store them onto \textit{Google App Engine Data Storage};
 	\item acquire the video, perform the beating plot, and to store them onto \textit{Google App Engine Data Storage};
 	\item get from the \textit{Google App Engine Data Storage} the electrovalves status set from the user through the Google Glass and to drive the electrovalves.
 \end{itemize} 
 
 The whole designed environment includes a program, written using the framework \textit{Qt}, for storing the recorded video from microscope.
 
 \section*{The Glasswear}
 \addcontentsline{toc}{section}{The Glasswear}
 The (Fig.\ref{Fig:GlasswearDiagram}) shows the structure of the Glasswear. From the Home Screen (Fig.\ref{Fig:GlasswearDiagram}a), using the voice trigger "\textit{Show Measurement}" or tapping on the "\textit{Measurement}" card (Fig.\ref{Fig:GlasswearDiagram}b) user is allowed to enter in the application (Fig.\ref{Fig:GlasswearDiagram}c). From this point, tapping and swiping, it's possible to navigate into the glasswear's menu (Fig.\ref{Fig:GlasswearDiagram}d-h) and choose which card has to be shown. \textit{View PH} (Fig.\ref{Fig:GlasswearDiagram}i) and \textit{View Temperature} (Fig.\ref{Fig:GlasswearDiagram}j) cards  plot on the card's left side the value of pH and temperature, respectively. While on the right side they show the average value. The microscope's video is shown by tapping on \textit{View Video} (Fig.\ref{Fig:GlasswearDiagram}k). The \textit{View Beating} card shows the graph of the beating associated to the video.
 \\
 
 \begin{figure}[h]
 	\centering
	\includegraphics[width=\textwidth]{Intro/block-app.eps}
 	\caption{Glasswear's Block Diagram}
 	\label{Fig:GlasswearDiagram}
 \end{figure}

 
 From the \textit{Drive Electrovalves} card (Fig.\ref{Fig:GlasswearDiagram}m), the user can set the value of each electrovalve. The main view of this card shows the status of each electrovalve (written in green if it is on and in red if it is off).
 \clearpage
 \begin{figure}[h]
 	\centering
 	\includegraphics[width=\textwidth]{Intro/electrovalves.eps}
 	\caption{Drive Electrovalves Steps}
 	\label{Fig:Electrovalves}
 \end{figure}
 
 The (Fig.\ref{Fig:Electrovalves}) shows the steps to toggle the status of the first electrovalve:
 \begin{enumerate}
 	\item (Fig.\ref{Fig:Electrovalves}a) shows the initial status of the whole electrovalves (all off);
 	\item tapping on the card and swiping the user is allowed to change the status of each electrovalve from the menu, as shown in (Fig.\ref{Fig:Electrovalves}b);
 	\item after that the electrovalve has been chosen, a toast message pops up (Fig.\ref{Fig:Electrovalves}c), and the new values of the electrovalves are shown. 
 \end{enumerate}
 
 To return on the main card of the glassware, the user has to tab on \textit{Back} item (Fig.\ref{Fig:Back}) from every menu.
 
 \begin{figure}[h]
 	\centering
 	\includegraphics[scale=.35]{Intro/back_menu}
 	\caption{\textit{Back} menu item}
 	\label{Fig:Back}
 \end{figure}
 
 To terminate the glassware, from every menu, the user has to swipe up to the final item and tab on \textit{Exit} item (Fig.\ref{Fig:Exit}).
 
 \begin{figure}[h]
 	\centering
 	\includegraphics[scale=.35]{Intro/exit}
 	\caption{\textit{Exit} menu item}
 	\label{Fig:Exit}
 \end{figure}
 
\section*{The Board}
\addcontentsline{toc}{section}{The Board}
\begin{figure}[h]
	\begin{center}
		\includegraphics[width=\textwidth]{thboard}
		\caption{The Board}
		\label{Fig:board}
	\end{center}
\end{figure}

The (Fig.\ref{Fig:board}) shows the top view of the system board. As can be seen it is composed by different interface and connections:
\begin{itemize}
	\item $5\ V$ power supply, required by the microcomputer on the Baglebone Black and by the conditiong circuits on the \textit{PCB};
	\item $24\ V$ power supply, required by the electrovalves;
	\item \textit{Ethernet}, to connect the board to the Internet;
	\item \textit{USB} connection, for the microscope;
	\item \textit{pH header}, to connect the pH sensor\footnote{\textbf{WARNING}: to ensure the correct functionality, user has to pay attention at this connection, since the pH sensor is a passive one, it has a polarity. The positive pin of the sensor has to be connected to the right terminal, looking fro the top (as shown in (Fig.\ref{Fig:board}))};
	\item \textit{temperature header}, to connect the temperature sensor;
	\item \textit{electrovalves header}, made by eight pairs of terminals, ordered as shown in (Fig.\ref{Fig:board}), from the top to the bottom.
\end{itemize}

The remaining two headers, shown in the top left corner of (Fig.\ref{Fig:board}) are thought for future application. In particular they are going to be useful for all those jobs that don't required the interaction with Google Glass, such as sensors calibration.

\section*{Video Storing}
\addcontentsline{toc}{section}{Video Storing}
The storing of the microscope video plays an important role of this system. It may be essential to review the recorded video  during the experiment and in order to fulfill this aim a \textit{Qt} program has been designed.\\
We chose \textit{Qt} because in this way the program is available for different operating systems, \textit{Linux}, \textit{Windows}, and \textit{MacOS}.\\

The program is very easy to use, the user just has to run the executable (Fig.\ref{Fig:icostoring}).

\begin{figure}[h]
	\begin{center}
		\includegraphics[scale=.5]{storing/Icon}
		\caption{Storing Microscope Video Program's Icon}
		\label{Fig:icostoring}
	\end{center}
\end{figure}

Once it has been launched, a console is opened (Fig.\ref{Fig:storingwindows}). The program checks every 20 seconds if a new video has been uploaded on the server. If so, the new video will be stored inside the computer (directory \texttt{C:/Video}) with the current data and hour as name in the following form: $YYYY-MM-DD-HH-mmss$, as shown in (Fig.\ref{Fig:stored}).\\

As shown in (Fig.\ref{Fig:storingwindows}) on the console the user can read all the information about what the program is doing.

\begin{figure}[h]
	\begin{center}
		\includegraphics[width=\textwidth]{storing/program-windows}
		\caption{Storing Microscope Video Console}
		\label{Fig:storingwindows}
	\end{center}
\end{figure}

\begin{figure}[h]
	\begin{center}
		\includegraphics[width=\textwidth]{storing/stored}
		\caption{Video Stored in the Folder}
		\label{Fig:stored}
	\end{center}
\end{figure}
 
 

\cleardoublepage
\part{Hardware}\label{HW}
%\chapter{Introduction}\label{ch:Iintroduction}
In this part of the thesis the hardware that has been designed in this system is explained. 
%It covers the sensors conditioning and the driver for the electrovalves as well as the \textit{DC-DC} converter which supplies the electrovalves when the $12\ V$ ones have to be used. has been request in order to use the system with two different kind of electrovalves, and which supplies them
The design specification for this part are:
\begin{itemize}
	\item make a sensor conditioning for pH and temperature sensors, see (Sec.\ref{phcon}) and (Sec.\ref{tempcon}) for more details;
	\item make a driver for electrovalves which must be able to drives two different kinds of electrovalves, both of them require $80\ mA$ but one type at $24\ V$ while the other one at $12\ V$. To fulfill this aim a \textit{DC-DC} converter has been designed, as well. See (Sec.\ref{sec:electrovalves}) for more details;
	\item make a \textit{PCB}  which supports and connects all the previous components and that is a \textit{capes} for the Beaglebone Black,  it has to be wedged on top of it, see (Chap.\ref{cha:PCB}) for more details.
\end{itemize}
\newpage

\chapter{Conditioning Circuit and Electrovalves Drivers}\label{ch:Analog}
\section{PH Conditioning}

\subsection{The Sensor}\label{sec:sensorPH}
A pH sensor is used to measure hydrogen ion activity 

\subsection{The Circuit}\label{sec:circuitPH}
As already explained in (Sec.\ref{sec:sensorPH}), the pH sensor is a passive sensor, which means no excitation source is required because the sensor itself generates its own electrical output signal. In particular any variation of pH in input is transduced in a voltage variation in output.\\

The pH sensor is also bipolar, this means the voltage output may be both positive and negative, as shown in (Fig.\ref{Fig:typicalph}).\\

\begin{figure}[h]
	\centering
	\includegraphics[scale = .65]{Sensors/tipicalph.eps}
	\caption{Typical pH-sensor transfer function}
	\label{Fig:typicalph}
	
\end{figure}

Resuming, it produce a voltage output that decreases linearly with pH of the solution being measured. The sensors give a sensitivity which ranges between $50$ and $70$ $mV/pH$ (it depends from sensor to sensor), this means that, in order to well observe this variation, an amplification stage may be required.

The (Fig.\ref{Fig:circuitph}) shows the adopted solution for conditioning the pH sensor. \\
Fist of all, since the pH sensor produces a bipolar signal and this application operates on a single voltage supply, the signal has been level shifted. To achieve this first challenge the operation amplifier \textit{U1} forces an off-set of $512\ mV$ to the pH sensor. Indeed, the \textit{LM4140A-1.0} is a high precision low noise \textit{LDO} (Low Drop Out) voltage reference which provides an accurate $1.024\ V$. This voltage has been halved by the $10\ K\Omega$ resistor divider. The \textit{U1} is in voltage follower configuration, thus its output should be equal to the input, and it biases the reference electrode of the pH sensor with $512\ mV$, at low impedance.\\
So, what the part of circuit made by \textit{U1} and \textit{LM4140A-1.0} does is to shift the bipolar pH sensor signal to an unipolar in order to be usable in the single-supply system.
 
\begin{figure}[h]
	\centering
	\includegraphics[width=\textwidth]{Sensors/phCircuit.eps}
	\caption{Conditioning circuit for pH sensor}
	\label{Fig:circuitph}
	
\end{figure}

Another challenge is given by the high impedance of the electrode. In fact the output impedance of the pH sensor is higher than $100\ M\Omega$. The circuit in (Fig.\ref{Fig:high}) shows a typical connection of this sensor where the output voltage is given by:

\begin{equation}
V_{out} \simeq V_{in} \text{ \textit{=} } V_{S} - I_{bias} \cdot R_{S}
\end{equation}

Thus, in order to reduce the error caused due to amplifier's input bias current a really low input bias current amplifier has to be chosen. For this reason, the 
\textit{LMP7721} is used, it is has an ultra-low input bias current ($3 \pm 17\ fA$).\\


\begin{figure}[h]
	\centering
	\includegraphics[]{Sensors/highimpedence.eps}
	\caption{Error caused by Amplifier's Input Bias Current}
	\label{Fig:high}
	
\end{figure}

In (Fig.\ref{Fig:circuitph}) both \textit{U1} an \textit{U2} are \textit{LMP7721}. The second amplifier with the \textit{DG412} represents a simple \textbf{PGA} (\textit{Programmable Gain Amplifier}), where the resistors have been chosen to give the following gains:
\begin{itemize}
	\item $1$, when $C_1$ is asserted and the others are denied;
	\item $2$, when $C_2$ is asserted and the others are denied;
	\item $4$, when $C_3$ is asserted and the others are denied.
\end{itemize}

The feedback capacitor is used to ensure stability and holds the output voltage during the switching times. Indeed, in these slice of time the output node would be floated without the capacitor.\\
This PGA stage introduces an additional offset error of  $75 \pm 470\ fV$, due to the bias current of the operational amplifier and $R_{ON}$ of \textit{DG412} ($25\ \varOmega$). That voltage combined with the \textit{LMP7721} offset (which is very higher than the first one) is approximately equal to $26\ \mu V$.\\

Since the environment in which the circuit is going to be used is a laboratory, so a really noisy place, the conditioning circuit for the sensor has to involve the design of a \textit{\textbf{low-pass filter}} in order to reject the noise.\\

The circuit shown in (Fig.\ref{Fig:filterCircuit}) is a third order low-pass filter with a \textit{Bessel} response. It has been designed in order to ensure a really flat-response in the pass band. \\
The parameter of the filter are:

\begin{enumerate}
	\item \textit{cutoff frequency} ($f_c$): $26.8\ Hz$, ;
\item \textit{stop band attenuation}: $-28.2\ dB$ at \textit{stop band frequency} ($f_s$) $60\ Hz$ (the line frequency in USA);
\item \textit{Quality factor} ($Q$): $0.65$;
\item \textit{filter order}: third;
\item \textit{filter response}: Bessel.
\end{enumerate}

It's important that $Q < 0.707$ because otherwise would be some peaking in the filter response. While, in this case, as shown in (Fig.\ref{Fig:filterbode}), roll-off at the cutoff frequency is greater.\\
This filter also behaves as \textit{anti-aliasing filter}, to prevent the aliasing components from being sampled during the analog to digital conversion.\\



\newpage
\clearpage

\begin{figure}[h]
	\begin{center}
	\includegraphics[height=.38\textheight]{Sensors/filterCircuit}
	\caption{Low-Pass Filter Schematic}
	\label{Fig:filterCircuit}

	\centering
	\includegraphics[height=.6\textheight]{Sensors/filterBode}
	\caption{Low-Pass Filter Frequency Response}
	\label{Fig:filterbode}
	\end{center}
\end{figure}
	
	\newpage
	\begin{figure}[t]
The output of this filter represent the input signal of the embedded ADC inside the Beaglebone Black.\\

Connecting together al this part we obtain the circuit in (Fig.\ref{Fig:pHConditioning}) where we can also see the general purpose input/output pin used to drive the \textit{PGA} and the analog input used for the pH sensor. As it is explained in (Chap.\ref{ch:firmware}) \texttt{AN\_2} is used to let the microcontroller know about the added offset.

\vspace{10mm}

\centering
\includegraphics[width = \textwidth]{pcb/ph_conditioning}
\caption{Acquisition Path for pH Sensor}
\label{Fig:pHConditioning}
\end{figure}






\newpage
\clearpage

\section{Temperature Conditioning}
\subsection{The Sensor}\label{sec:sensorT}

\subsection{The Circuit}

As already explained in (Sec.\ref{sec:sensorT}), the temperature sensor is a active sensor, which means excitation source is required because the sensor is resistor based so a current must be passed through it. Then, the corresponding voltage has to be measured in order to determine the temperature value.\\
So, any variation of temperature in input is transduced in a resistance variation in output.\\

The (Fig.\ref{Fig:temperatureCircuit}) shows the adopted solution for conditioning the temperature sensor.\\
In this connection the temperature sensor is excited by $680\ \mu A$, so the voltage $V_T$ is given by the following equation:

\begin{equation}
V_T \text{ \textit{=} } 680 \mu \cdot R_T
\end{equation}

\begin{figure}[h]
	\begin{center}
		\includegraphics[height=.38\textheight]{Sensors/temperatureCircuit}
		\caption{Conditioning Circuit for Temperature Sensor}
		\label{Fig:temperatureCircuit}
	\end{center}
\end{figure}

In order to provide this amount of current a \textit{LT3092} is used. It can supply an output current equal to:

\begin{equation}
I_{OUT} \text{\textit{ = }} 10\mu \cdot \frac{R_{SET}}{R_{OUT}}
\end{equation}

To ensure the stability of the component, a feedback capacitor of $1\ \mu F$ is exploited. The transistor $Q_1$ is used to avoid the self-heating of the temperature sensor: when the temperature value has to be sampled $In$ is denied, for the remaining time $In$ is asserted, in this way the resistive sensor is by-passed, and the \textit{Joule} effect is avoided.\\

For the same reason exposed in (Sec.\ref{sec:circuitPH}), also in this case a filter is needed, and, since the voltage value is almost the same, the used filter is equivalent to the previous on.\\

\begin{figure}[h]
	\centering
	\includegraphics[width = \textwidth]{pcb/temperature_conditioning}
	\caption{Acquisition Path for pH Sensor}
	\label{Fig:temperatureConditioning}
\end{figure}

 
\section{Electrovalves Driver} \label{sec:electrovalves}

In order to allow the reverse control, from Google Glass to electrovalves, it is important to design a circuit that has to drive them from digital values provided by the \textit{Beaglebone Black} ($0$ equal to $0\ V$, $1$ equal to $3.3\ V$ and in both cases the maximum suppliable current is only $6\ mA$).\\

The electrovalves used are \href{http://www.festo.com/net/SupportPortal/Files/10026/MH1_VO_ENUS.pdf}{FESTO solenoid valves MH1}. They require a voltage of $24\ V$ in DC and a current of $80\ mA$.\\

\begin{figure}[h]
	\centering
	\includegraphics[width = .6\textwidth]{Driver/electrovalves}
	\caption{Electrovalves}
	\label{Fig:EV}
\end{figure}

To fulfill this aim a low-side switch MOS has been used, as shown in (Fig.\ref{Fig:driverEV}).

\begin{figure}[h]
	\centering
	\includegraphics[]{Driver/driverEV}
	\caption{Driver for electrovalves}
	\label{Fig:driverEV}
\end{figure}

The \textit{Fairchild Semiconductor BS170 N-Chnnel MOS} has been chosen because of its low price and its capability of supporting a drain-source voltage of up to $60\ V$ and supplying a continuous drain current of up to $500\ mA$.\\

To protect the \textit{BS170} from reverse inductive current surges due to the solenoid of the electrovalve, a \textit{Vishay Semiconductors 1N4148 Diode} is used. It is able to support  a reverse voltage of up to $75\ V$ and a continuous forward current of up to $150\ mA$.\\

As shown in (Fig.\ref{Fig:EV}) the number of electrovalves used is eight, so the previous driver has been replicated in order to obtain the circuit in (Fig.)

\begin{figure}[h]
	\centering
	\includegraphics[width = \textwidth]{Driver/electrovalves_circuit}
	\caption{Electrovalves Driver Circuit}
	\label{Fig:electrovalves_circuit}
\end{figure}

\section{DC-DC Converter}

Since this system is supposed to drive different kind of electrovalves, which require a different value of voltage (but not in current) a \textit{DC-DC converter} is used in order to convert the voltage supply from $24\ V$ to $12\ V$. \\

Using the \textit{LT1374} the circuit in (Fig.\ref{Fig:DCDC}) has been designed, it is a constant frequency (equal to $500\ Hz$), current mode buck converter. It has an embedded clock and two feedback loops to control the duty cycle of power switch.\\
%This converter is able to supply a current of up to $2.5\ A$, that is definitely  bigger than the maximum required current, in which all the electrovalves are on in the same time: $80 \cdot 8 \text{ = } 640\ mA$.

\begin{figure}[h]
	\begin{center}
		\includegraphics[width = \textwidth]{Sensors/dcdcCircuit}
		\caption{DC-DC circuit}
		\label{Fig:DCDC}
	\end{center}
\end{figure}

%The output capacitor not only behaviors like a low-pass filter in order to reduce the output spike, but it is used also to maintained the efficiency of \textit{DC-DC} over the output current range 

Trough a low-side switch, made with the \textit{BS170}, it is possible to shutdown the converter from the software. When the \textit{LT1374} is in shutdown, the supply current is reduced to $20\ \mu A$. As it is explained in (Chap.\ref{ch:firmware}) this is used to prevent regulator from operating when $24\ V$ are required.\\

\subsection{Components Choice}

All the components have been chosen with attention and for some reasons that are going to be explained belong.
\subsubsection{FEEDBACK RESISTORS}
The main behavior of the feedback pin on the \textit{LT1374} is to set the output voltage, and this deals with selecting the resistors $R_1$ and $R_2$. They are related from each others by the following equation:
\begin{equation}
R_1 = \frac{R_2 \cdot \left( V_{OUT} - 2.42\right)}{2.42}
\label{Eq:feedback}
\end{equation}
As suggested on datasheet of the component, the resistor between feedback pin and ground is $4.99\ k\varOmega$. So, from the (Eq.\ref{Eq:feedback}) results that the value of $R_1$ is $19.6\ k\varOmega$.
\subsubsection{INDUCTOR}
The choice of inductor is a trade-off among:
\begin{itemize}
	\item \textit{physical area}, lower values of inductor mean lower size;
	\item \textit{output current}, higher values of inductor allow more output current because they reduce peak current ($I_{SW\left(PEAK\right)} \propto 1/L$)
	\item \textit{ripple voltage}, higher values of inductor reduce the output ripple voltage.
\end{itemize}
A good choice is represented by $10\ \mu H$, with this inductor the maximum current peak (Eq.\ref{eq:peak}) is equal to $1.24\ A$.
\begin{equation}
I_{SW\left(PEAK\right)} \text{ = } I_{OUT} + \frac{V_{OUT \cdot \left(V_{IN} - V_{OUT}\right)}}{2 \cdot f \cdot L \cdot V_{IN}}
\label{eq:peak}
\end{equation}
\subsubsection{OUTPUT CAPACITOR}
The output capacitor determines the output ripple voltage, for this reason a small \textit{Effective Series Resistance} (ESR) is required.\\
The frequency operation of \textit{LT1374}, as already said, is equal to $500\ Hz$ and at this frequency any polarized capacitor is essentially resistive. As suggested from datasheet, for typical \textit{LT1374} application the ESR has to range from $0.05\ \varOmega$ to $0.2\ \varOmega$, for this reason the output capacitor is a \textit{solid tantalum capacitor}. The choice of $100\ \mu F$ is a good trade-off between output ripple voltage and physical area.
\subsubsection{SCHOTTKY DIODE}
The chosen diode is \textit{On Semiconductor MBR330} because of its capability of supporting a $3\ A$ average forward current and $30\ V$ reverse voltage.\\
Indeed, the reverse voltage is approximately $12\ V$ (the output voltage), while
the average forward current is given by the (Eq.\ref{Eq:avgCurrent}).
\begin{equation}
	I_{D \left(AVG\right)} \text{ = } \frac{I_{OUT \cdot \left(V_{IN} - V_{OUT}\right)}}{V_{IN}}
	\label{Eq:avgCurrent}
\end{equation}
The (Eq.\ref{Eq:avgCurrent}) will never yield values higher than $3\ A$, neither in worst-case scenario.
% represented by the overloaded (not the shorted) output. In this case, in fact, the output current of the \textit{LT1374} increase to a value of $5.7\ A$ and the average forward current 
\subsubsection{BOOST CAPACITOR}
The boost capacitor has been chosen based on the voltage that has to support, which is basically equal to the output voltage ($12\ V$) and the (Eq.\ref{Eq:minCap}) provided by \textit{LT} on datasheet of the component. In this application result that its minimum value is equal to $1.5\ nF$.

\begin{equation}
C_{MIN} \text{ = } \frac{\left(I_{OUT} / 50 \right) \cdot \left( V_{OUT} / V_{IN}\right)}{f \cdot \left(V_{OUT} - 3\right)}
\label{Eq:minCap}
\end{equation} 

\subsection{Relay}

Since the electrovalves may need $12\ V$ or $24\ V$ a way to switch between this two voltages supplies is needed. The circuit shown in (Fig.\ref{Fig:relay}) has been used for this purpose.\\
It allows the switching trough the firmware.The circuit uses a optocoupler, the \href{http://www.keepjump.com.tw/DataSheet/Others/DPC-817C.pdf}{DPC-817C}, to isolate the Beaglebone Black to the relay, and prevent in this way that pounces produced by magnetic part of relay reach the general purpose pin of the Beglebone itself. The relay used is the \href{http://www.rlocman.ru/i/File/dat/Omron/Relays/G5LA145DC.pdf}{G5LA1CF24DC} and, as happened in (Sec.\ref{sec:electrovalves}), a diode has been exploited to preserve the transistor used as voltage controlled switch from broking. 

\begin{figure}[h]
	\begin{center}
		\includegraphics[width = .9\textwidth]{Sensors/relay}
		\caption{Relay circuit}
		\label{Fig:relay}
	\end{center}
\end{figure}

Summary, in the circuit represented  in (Fig.\ref{Fig:relay}) the \textit{Supply} is the voltage which is going to be provided to the electrovalves. When the \textit{In} signal is asserted, the relay is in its normal connection, son \textit{Supply} is tied to $24\ V$. On the other hand, when \textit{In} is denied relay is excited and \textit{Supply} is tied to $12\ V$.
\newpage

\begin{figure}[t]
	\centering
	\includegraphics[width = \textwidth]{pcb/supply}
	\caption{Electrovalves Supply Circuit}
	\label{Fig:supply}
\end{figure}

The circuit in (Fig.\ref{Fig:supply}) shows the connection between \textit{DC-DC} circuit and the switching one, all of this is necessary to correctly supply the different kind of electrovalves.



\chapter{PCB}\label{cha:PCB}
\begin{figure}[h]
	\centering
	\includegraphics[width = \textwidth]{circuitfinisched}
	\caption{Final system mounted on a PCB}
	\label{Fig:finalcircuit}
\end{figure}


All the hardware and the circuit described so far has been mounted on a Printed Circuit Board (\textit{PCB}) that is supposed to be a \textit{cape} of the Beaglebone Black. This means that the whole physical structure has been made to be attached on the headers of the board. In (Fig.\ref{Fig:circuit}) all the circuit that has to be made on the \textit{PCB} is shown.\\

The PCB has been designed in a two-layer board ($100\ x\ 100\ mm$) using the \textit{National Instruments Circuit Design Suite},  in particular \textit{Multisim} to carry out the schematic and \textit{Ultiboard} for what concern the PCB.\\
The design has followed the guidelines for reduced electromagnetic interface (\textit{EMI}). First of all, since Surface-mount devices (\textit{SMD}) are better than Through-hole components (\textit{THD}) in dealing with RF energy, because of reduced inductances and closer component placements available (\cite{PDGFRE}), where there was a choice, \textit{SMD} components have been used. \\

Particular attention has been paid for the \textit{DC-DC} converter layout, because a wrong \textit{PCB} design for switching power supply often means failure. Moreover, in these terms, what is good for \textit{EMI} is also good in terms of functional stability for the regulator.\\




\begin{figure}[t]
	\begin{center}
	\includegraphics[angle = 90, width = \textwidth]{pcb/circuit1}
	\caption{Schematic of Complete Circuit}
	\label{Fig:circuit}
	\end{center}
	
	
\end{figure}

 \clearpage
 
 The circuit in (Fig.\ref{Fig:High}) is schematically a common buck regulator. In that figure, with a blue circle, is highlighted the high speed switching current path: the loop which produces the highest \textit{EMI}. Indeed, in  the blue loop flows a fully switched alternated current, and for this reason it is also referred as \textbf{hot loop}.

\begin{figure}[h]
	\centering
	\includegraphics[width = \textwidth]{pcb/dcdcloop}
	\caption{\textit{DC-DC} High Speed Switching Path}
	\label{Fig:High}
\end{figure}

In order to ensure clean switching and reduce \textit{EMI} the minimum lead length is required for reducing the radiating effect of the hot loop as much as possible: and so it has been done, as can be seen from (Fig.\ref{Fig:dcdc-pcb}).

\begin{figure}[h]
	\centering
	\includegraphics[width = \textwidth]{pcb/dcdc-pcb1}
	\caption{\textit{PCB} Layout of \textit{DC-DC}}
	\label{Fig:dcdc-pcb}
\end{figure}

In (Fig.\ref{Fig:dcdc-pcb}) the \textit{PCB} layout for \textit{DC-DC} converter has been shown, in which we can see that: 
\begin{itemize}
	\item the magnetic radiation is minimized by keeping by keeping catch diode and the input capacitor leads as short as possible;
	\item the electric radiation is minimized by reducing the area and length of all traces connected to the switch and boost pins.
\end{itemize}

Moreover, in order to reduce the noise on the feedback, that is translated in error on output voltage, the switch node and the feedback resistors are kept as far as possible from each others.\\

The (Fig.\ref{fig:PCB}) shows the \textit{PCB} layout without ground plane (Fig.\ref{subfig-1:pcb}), then with both top and bottom ground plane (Fig.\ref{subfig-2:pcbboth}) which help with heat dissipation and \textit{EMI} reduction. In (Fig.\ref{subfig-2:pcb3d}) and (Fig.\ref{subfig-2:pcbphoto}) the 3-D model generated using \textit{Ultiboard} and the final real result are compared.


\begin{figure}[h]
	\subfloat[\textit{PCB} Layout Withot Ground Plane\label{subfig-1:pcb}]{%
		\includegraphics[width=.5\textwidth]{pcb/PCB}
	}
	\subfloat[\textit{PCB} Layout With Both Top and Bottom Ground Plane\label{subfig-2:pcbboth}]{%
		\includegraphics[width=.5\textwidth]{pcb/PCB-both}
	}\\
	\subfloat[\textit{PCB} 3-D Model View\label{subfig-2:pcb3d}]{%
		\includegraphics[width=.5\textwidth]{pcb/pcb-up}
	}
	\subfloat[\textit{PCB} Real Result\label{subfig-2:pcbphoto}]{%
		\includegraphics[width=.5\textwidth]{pcb/foto}
	}
	\caption{PCB of the System}\label{fig:PCB}
\end{figure}


\part{Firmware}
\chapter{Introduction}\label{ch:IIintroduction}
\begin{figure}[h]
	\centering
	\includegraphics[scale = .9]{bo}
	\caption{High Level System's Block Diagram}
	\label{Fig:logicalview}
	
\end{figure}

The term \textit{Embedded Systems} is widely used nowadays, because its definition is  very generic: any device that includes a programmable computer, but is not itself a general-purpose computer is referred as Embedded Systems.

Indeed this definition covers a huge number of systems, from  a vending machine to a smartphone, form a electric toothbrush to a body computer of a car.\\
Peculiarity of an Embedded System is that it has hardware and software parts, takes advantage of application characteristics to optimize the design, interacts with the external environment using sensors and actuators.

\textbf{Embedded Linux} systems are a subset of embedded systems that are enhanced by a Linux operating system (\textit{OS}), here the integration of high-level Linux software and low-level electronics represents a paradigm shift in embedded systems development \cite{EBB}. It allows an easy way to design systems that can meet future challenges in smart buildings, the Internet of Things (\textit{IoT}), human-computer interaction(\textit{HCI}), robotics, and many other applications.

In this thesis project, as can be seen from (Fig.\ref{Fig:logicalview}), the system may be intended as an \textit{IoT} one. Indeed, in order to allow the interfacing between sensors and electrovalves with the Google Glass, I had to use an Internet connection, represented by the Google App Engine. And, while the Google Glass can interact directly with this server, the hardware described in (Part.\ref{HW}) needs a micro computer to be connected with that server.

\section{A Brief Introduction To IoT} 

The term Internet of Things is broadly used to describe the extension of the web and the Internet for the physical objects or "\textit{things}". It is important to realize that the physical connection of the human Internet and the ones of the \textit{IoT} are the same. And even if the \textit{IoT} term has be coined recently, the majority of web traffic nowadays is non-human \cite{ALTIOT}.

Indeed, the massive usage of online services, in part due to the smartphone revolution, has increased the interactions between servers, machine-to-machine (\textit{M2M}) communications. This trend can only increase since it is expected an explosion of wearable systems, such as smartwatch, smartglass and so on..



\begin{figure}[h]
	\centering
	\includegraphics[scale = .6]{traffic}
	\caption{Distribution of Internet of Things}
	\label{Fig:traffic}
	
\end{figure}

In the near future, \textit{Cisco IBSG} predicts there will be $25$ billion devices connected to the Internet by $2015$ and this number is destinate to reach $50$ billion by $2020$, see (Fig.\ref{Fig:cisco})\cite{CISCOIOT}. It is also important to see that these number are based on the grown of $2011$ and do  not consider the rapid advances in device technology. So the number shown in (Fig.\ref{Fig:cisco}) have to be considered as minimum.

\begin{figure}[h]
	\centering
	\includegraphics[scale = .65]{cisco}
	\caption{The \textit{IoT} Evolution Prediction}
	\label{Fig:cisco}
	
\end{figure}


The IoT concept, that has been exploited even in this thesis, is that if physical sensors and actuators can be linked to the Internet , then a lot of new applications and services are possible \cite{EBB}.

\section{The Beaglebone Black}

For this thesis, the Embedded Linux system has been developed using the \textit{Beaglebone Black}.\\

\begin{figure}[h]
	\centering
	\includegraphics[scale=.55]{beaglebone}
	\caption{The \textit{Beaglebone Black}}
	\label{Fig:beaglebone}
	
\end{figure}

The \textit{Beaglebone Black} is a low-cost, fully open-source (that means everyone can downloads and uses the hardware schematics and layout in order to build up a complete product design), single board micro computer (\textit{SBC}) powered by Linux. 

There are a lot of \textit{SBC} on the market, I chose the Beaglebone Black for the following reasons:
\begin{itemize}
	\item low price, it costs $\$50 $;
	\item powerful, the \textit{Sitara\texttrademark AM335x ARM\textregistered Cortex\texttrademark-A8} processor from Texas Instruments (\textit{TI}) can run up to $1\ GHz$, allowing up to $2$ billion instructions per second;
	\item $4GB$ on-board embedded multi-media card (\textit{eMMC}), this means the Beaglebone can boot without a \textit{SD} card, unlike other boards like Raspberry PI. It allows faster boot, approximately $10$ seconds;
	\item $512MB$ \textit{DDR3} of system memory;
	\item Ethernet processor, which supports \textit{DHCP}, so it can be directly conncted to a network
	\item \textbf{expansion headers}, $92$ pins with $65$ \textit{GPIOs}, $8$ analog output, $7$ analog input, $3$ different voltage supplies ($5\ V$, $3.3\ V$, and $1.8\ V$), and the whole most used digital interface. All these make the Beaglebone Black built to be interfaced to. 
	
\end{itemize}
Over these, the board has other important feature that are not used in this thesis, but may be used in the future to enhance it. An example is given by the 2 Programmable Real-time Units (\textit{PRU}) that, unlike the definitions of Embedded Linux, make this board usable for real-time applications.

As already said in (Chap.\ref{ch:Analog}), this board has been connected to a custom \textit{capes}, a daughterboard that can be attached to the two headers on the Beaglebone itself.


\chapter{Embedded Linux inside The project}\label{ch:firmware}
In this section the different tasks that Beaglebone has o perform are described in detail. In this project the \textit{Linux distro} used is \textit{Debian}. It has been chosen among the wide possibilities because ensure a good level of stability, and moreover it has a very good repository. 

\begin{figure}[h]
	\centering
	\begin{tikzpicture}[
	level 1/.style={sibling distance=40mm},
	edge from parent/.style={->,draw, black},
	>=latex, scale=0.75, every node/.style={scale=0.70}]
	
	% root of the the initial tree, level 1
	\node[root] {Embedded Linux Firmware}
	% The first level, as children of the initial tree
	child {node[level 2] (c1) {Video Acquisition}
		child{node[level 3] {Beating Plot Computation}}
		%child{node[level 3] {Electrovalves Driver}}
		%child{node[level 3] {Power Supply}}
		}
	child {node[level 2] (c2) {Sensors Data Acquisition}}
	child {node[level 2] (c3) {Electrovalves Driver}
		};

	
	\end{tikzpicture}
	\caption{Embedded Linux Firmware - Blocks Subdivision}
	\label{Fig:firmware}
\end{figure}

In (Fig.\ref{Fig:firmware}) are shown the three principle subsection of the firmware. For each of them a \textit{C++} program has been written, and their executable files are going to be run through a \textit{bash} script, which is supposed to be launched at the system boot. The (List.\ref{code:recordVideo}) shown this bash script that behaves as shown in (Fig.\ref{Fig:bash}).

First of all the \textit{bash} script loads the \textit{Device Tree Overlay}. The pins of \textit{Sitara\texttrademark} microprocessor are \textit{multiplexed}, it means that each pin of the two header may assume different behavior (\textit{GPIO}, \textit{SPI}, \textit{I2C}, and so on..). During the booting a default value for this mux is loaded, and for certain pins, this default value differs from the actual value that the system needs. In order to change those values I wrote a Device Tree Overlays (\textit{DTO}) which has to run at the beginning, immediately after the booting, in order to explain how is it work the Flattened Device Tree (\textit{FDT}) has to be introduced.

The last releases of Linux kernels for ARM boards use \textit{FDT}, it is a data structure that describes the hardware on board. Thus, it describes every component presents on the board, from the user \textit{LEDs} to the \textit{CPU}. Using this \textit{FDT} it is possible to write a \textit{device tree overlay} in order to change to mode of use of each pin. The (List.\ref{code:pins}) changes the status of $10$ pins from their default value, in particular sets the pins used to drive the valves, to allow the current on resistor sensor, and to change the power supply voltage on the valves to be digital output pins, with a pulldown resistor.

To record the video from the microscope I used the program \textit{\textbf{Video for Linux v.2}} (\textit{V4L2}), before start record the \textit{bash} script sets the video format as \textit{MJPG} and the dimension as $320\ x\ 240\ @ 30\ fps$. After that it runs in background the program to update the electrovalves status and to acquire the sensors value. As it is explained in (Sec.\ref{sec:valves}), these programs run out of the infinite loop because they have an infinite loop inside.

Inside the infinite loop the \textit{bash} script launch the recording of $10$ seconds microscope video, compresses the latter and elaborate it in order to compute an image with the beating plot. Finally, it uploads (using the library \textit{CURL}) video and image and update the flags that in order to indicate to Google Glass App and Video Storing software that a new video has been updated.

\begin{figure}[h]
	\centering
	\includegraphics[height = .74\textheight]{firmware/bash}
	\caption{Bash Script Flow Chart}
	\label{Fig:bash}
	
\end{figure}
\clearpage

\section{Video Processing and Beating Plot}

As already introduced once the microscope video has been recorded in a \textit{MP4} format it has to be processed in order to compute the beating rate. The code described in this section is shown in (List.\ref{code:beating}). 

To achieve this goal, the \textit{OpenCV} library has been used. It is a specific library for computer vision.

As can be seen the program captures the video from the file and, elaborating frame by frame it computes the point to be plotted. In particular it sores the first frame of video, then each point is given by the mean different color density between the under study frame and the first one.

The results has been stored in file called \textit{video\_data.dat}. This file is going to be the input for the \textit{gnuplot} task that is in charge to plot the beating data.

\section{Sensor Data Acquisition}

From the main \textit{bash} script another one is launched for acquiring the sensors data. This script is shown in (list.\ref{code:sensorbash}) and as can be seen it is an infinite loop where two operations are performed:
\begin{enumerate}
	\item cleaning the previous data from the datastore of Google App Engine server, otherwise plots on the Google Glass are impossible to be viewed;
	\item acquiring and uploading $100$ new sensors data. 
\end{enumerate} 

This second step is performed by another task shown in (List.\ref{code:sensors}), in which three objects have been used (Fig.\ref{Fig:umlsensors}). 

How it works: using a for loop, program runs $100$ times the reading of pH and temperature values, after that the reading step has been accomplished, the program is in charge to store these value on the Datastore of Google App Engine then it waits for $100\ ms$. Two important features has to be highlighted:
\begin{enumerate}
	\item as already introduced in (Sec.\ref{tempcon}), there is a \textit{BJT} inside the temperature conditioning circuit that acts as software-controlled switch. This is necessary to avoid the self-heating of temperature sensor and ensure a correct measure. This software-controlling may be seen from the code, indeed there is a digital output called \textit{temperature\_disable} that is normally denied, and it is asserted only when a temperature measure has to be carried out;
	\item as already introduced in (Sec.\ref{phcon}), the pH measuring process need to know exactly how much is the amount of added offset. Indeed, in order to make the pH sensor response unipolar, the circuit polarizes the negative pin of pH sensor at approximately $512\ mV$. Instead of using the theoretical value it is better to measure it just before every pH measure process. This is easily viewable from the code. 
\end{enumerate}   

\begin{figure}[h]
	\centering
	\includegraphics[width = \textwidth]{Sensors/class}
	\caption{Sensor UML Description}
	\label{Fig:umlsensors}
	
\end{figure}

\section{Electrovalves Updating Task}\label{sec:valves}

The code of task which performs the driving of the valves is shown in (List.\ref{code:Electrovalves}). Also this program uses the \textit{curl} library to communicate with the Google App Engine, and also this program has been built with an infinite loop.

What the task does is to download from the Google App Engine the status of each electrovalve that has been set by the user through the Google Glass. These status value are then stored onto a vector, to be written in a second step onto the output pins. As explained in (Sec.\ref{sec:electrovalves}) each valve is connected to a \textit{MOS} transistor that acts as voltage-controlled switch. The voltage which drives the transistor is the digital value written onto the output pins.

The task may be called with a parameter, the suply voltage value for valves:
\begin{verbatim}
./ValvesUpdating [12 | 24]
\end{verbatim}
If the parameter is missing, $24\ V$ as voltage are assumed as default. This parameter is going to drive both the \textit{DC-DC} enable and relay selector. In fact, if the voltage is the default one, the \textit{DC-DC} does not need to work, and so, in order to reduce the power consumption, be disabled.
\cleardoublepage

\part{Software}
Rising in the logical level we find the \textit{software}. In this part of thesis the python app which controls the server and the video storing program are going to be explained.

The first plays the important role to link the circuit, the embedded Linux system, with the Google Glass App, whenever the user is. While the second is in charge to store every new video that has been recorded, inside an hard drive of a generic computer, in order to be watched in a second time.  

More details of each of two are shown in the following sections.
\chapter{Google App Engine}
\begin{figure}[h]
	\centering
	\includegraphics[scale=.20]{applogo}
	\caption{Google App Engine Logo}
	\label{Fig:applogo}	
\end{figure}

Google App Engine is a Platform as a Service (\textit{PaaS}) which allows users to build and run web applications on Google  infrastructure.

Thus, when a programmer writes a web application that runs on App Engine, the software is going to run on the Google servers, somewhere in the Google cloude \cite{UGAE}. This solution is very useful because allows everyone to write the own web program without purchasing and maintaining the needed hardware and network.

For the time being Google App Engine supports web applications written in a variety of the main programming languages for the web: \textit{Java}, \textit{Python}, \textit{PHP}, and \textit{GO}. Since the App Engine was first built in $2008$ for Python, this is the language used to carry out the web application for this thesis.

Over the possibility to create our own web application without dealing with hardware, the Google App Engine has been chosen because of its datastore and memcache, which allow a simple memorization of data, automatic scaling and load balancing, and more others important features that are going to be shown below.

To create a more scalable and hierarchical web application the \textit{Flask framework} has been used. Flask is a lightweight web application framework written in Python. It is based on the \textit{WSGI} toolkit and Jinja2 template engine and it is distributed with a \textit{BSD} license.\\ Flask provides an easy way to organize a web application which allows to build up complex website easy to manage. It has no database abstraction layer, indeed for this purpose, in this application, it supports the Google App Engine.

\begin{figure}[h]
	\centering
	\includegraphics[scale=.8]{flask}
	\caption{Flask Logo}
	\label{Fig:flasklogo}	
\end{figure}

Before to explain how the code is made, the App Engine Datastore and Memcache are going to be introduced.\\
Using  the Google App Engine means do not have access to traditional database like Oracle and MySQL. In fact, App Engine uses \textit{\textbf{Google Datastore}}, which is easier to use because it takes more of a hierarchical object-oriented storage approach. That approach allows to ensure efficient application scalability.
Thus, the Datastore holds data objects known as \textit{entities}. An entity is composed by one or more \textit{properties}. Making a comparison with the object oriented it is possible to say that the properties are the fields of objects. So, likely the field, a property may support different data types such as integer, float, string and so on... Any entity holds its \textit{kind} and \textit{key}, which categorize and uniquely identify it, respectively. 

The \textit{memcache} results very useful to speed up common Datastore queries, indeed inside the Google Datastore, distributed memory cache are used for storage.
For this reason memcache is used to store the status of the electrovalves, in fact in this way the delay caused by the server is highly reduced.

\section{The web application}

The main code which runs on server is shown in (List.\ref{code:server}). It uses two models to describe the data supposed to be stored:
\begin{enumerate}
	\item \textit{Sensor}, contains the list of sensors used from the system, through the method \textit{sensor\_list} it returns that list;
	\item \textit{DataPoint}, contains all the points that have to be plotted, the field value represent the $y-axes$ while time the $x-axes$. While the first field has to be passed at the moment in which the class is constructed, the \textit{date} field is auto-filled with the current data. This class has two public methods to interact with: 
	\begin{itemize}
		\item \textit{point\_for\_sensor}, that with the name of sensor as parameter returns the list of \textit{DataPoint} associated;
		\item \textit{oldest\_point}, that returns the last \textit{DataPoint} inserted.
	\end{itemize} 
 \end{enumerate} 
 
 Below the different services offered by this web application are listed:
 \begin{itemize}
 	\item \textit{process\_values}, available at the path \url{/sensor_values}. User can access to this service using both \textit{HTTP} \textit{GET} and \textit{POST} methods. With the first one, it  returns the \textit{JSON} (\textit{JavaScript Object Notation}) containing all the data regarding the whole sensors stored inside the Datastore, in this format: \{ value, timestamp, sensor\_name \}. On the other hand, when the used method is \textit{POST} the application takes from the POST parameters those named \textit{"sensor"}, indicating the sensor's name, and \textit{"value"}, which is the sensor's value that has to be plotted. Thus, in this last way what the application does is to add a new data point for a sensor and store it inside the Datastore.
 	\item \textit{print\_names}, available at the path \url{/sensor_names}. This function is accessible only through HTTP \textit{GET} method and returns a \textit{JSON} object with the content of sensor's name list.
 	\item \textit{graphing\_data}, available at the path \url{/graphing_data}. To access to  this function a HTTP \textit{GET} method is required. This \textit{GET} request must contain two parameters: \textit{"sensor"}, the sensor name, and \textit{"first\_timestamp"}, which is the timestamp of the first point that has to be plotted. This last parameter is not mandatory and if it is missed it is assumed equal to $0$. What this function does is to return a \textit{JSON} object containing the whole DataPoint for the given sensor that have timestamp higher than \textit{"first\_timestamp"}.
 	\item \textit{clear\_data}, available at the path \url{/clear}. This function is in charge to deletes all the points for each sensor.
 	\item \textit{set\_picture}, available at path \url{/picture/submit}. This function is accessible with both HTTP \textit{POST} and \textit{GET} methods, and in both cases its behavior is to load a new image, containing the beating plot, onto Datastore. In the first case the image has be passed ad \textit{POST} data, and named as \textit{"Image.jpg"}. While th second case is supposed to be a browser way to upload the image. In fact, browsing to that \textit{URL} what the user sees is shown in (Fig.\ref{Fig:http}), a view where it is possible to select the picture file and submit it.
 	
 	\begin{figure}[h]
 		\centering
 		\includegraphics[scale=.8]{http}
 		\caption{Picture Submitting from the Browser}
 		\label{Fig:http}	
 	\end{figure}
 	
 	\item \textit{view\_picture}, available at the path \url{/picture/view}. Using this function, with a HTTP \textit{GET} method causes the access to the image previously stored inside the Datastore.
 	\item \textit{set\_video}, available at the path \url{/video/submit}. As happened to the beating plot image, this function is accessible with both HTTP \textit{POST} and \textit{GET} methods. And, as before, it is used to load onto Datastore the microscope video. Using the \textit{POST} method means storing the data passed as parameter and named \textit{"Video.mp4"}. On the other hand, in case of GET method what the user sees on the browser is shown in (Fig.\ref{Fig:http1}). As before, through the two buttons, user can choose the video to store and submit the request.
 	
 	\begin{figure}[h]
 		\centering
 		\includegraphics[scale=.8]{http1}
 		\caption{Video Submitting from the Browser}
 		\label{Fig:http1}	
 	\end{figure}
 	
 	\item \textit{view\_video}, available at the path \url{/video/view}. This function, accessible with HTTP \textit{GET} method, accesses to the microscope's video on the Datastore and returns it.
 	
 	\item  \textit{add\_electrovalve}, available at the path \url{/add/electrovalve}. This function is in charge to store the valve's status inside the memcache. It is accessible through HTTP \textit{POST} method, which must have these two following parameters:\textit{ "name"}, that uniquely identifies the valves (it is \textit{"EV"} follows by the number of the valve, for example the first one is \textit{"EV1"}), and \textit{"status"}, which indicates the actual status of the valve (it may be \textit{on} or \textit{off}).
 	
 	\item \textit{get\_electrovalve}, available at the path \url{/electrovalves/<name>}. It is accessible through HTTP \textit{GET} method and returns the value of the valve identified with \textit{<name>} field inside the \textit{URL} (for example, if someone wants to read the status of the first valve, has to use the path \url{/electrovalves/EV1}).
 	
 	
 	
 \end{itemize}
 
\chapter{Microscope Video Storing}
An important role of an experiment in organ-on-a-chip applications, as well as the whole biomedical field, is given by the video analysis. In fact, analyze the video just one time in real-time, during the experiment is running is never sufficient. So, what the system described in this thesis needs to be usable, is a way to memorize the video in order to be used, watched, in a second time.

To fulfill this aim I designed a console application using the \textit{Qt} framework. The motivation that brought me to use this framework is that the same code can run in different platform. In other words, this application can run on \textit{Linux}, \textit{Windows}, and \textit{MacOS}  indistinctly. This is a big advantage, because in a laboratory environment there are many researchers, and it is very easy to encounter different operating systems.

 In (App.\ref{code:video}) the listings of this application are shown. As can be seen, the source codes are divided in such a way to ensure high hierarchical efficiency.
 
 Indeed, the main (List.\ref{code:mainstoring}) of this application just instantiates a \textit{MyTimer} object. This \textit{MyTimer} object (List.\ref{code:timer}) is in charge to generate an interrupt every $10\ sec$, and it starts from its creation. When this interrupt comes, the timer uses the \textit{Downloader} object to check if  a new microscope video has been uploaded, and if so, download it and reset the flag that points up the new video status. The code of \textit{Downloader} is shown in (List.\ref{code:downloader}). The videos are stored inside the directory \textit{C:/Video} and their names correspond to the date and hour of download. 
 
 \begin{figure}[h]
 	\centering
 	\includegraphics[width=\textwidth]{storing/videostoring}
 	\caption{Video Storing: Block Diagram}
 	\label{Fig:videostoring}
 	
 \end{figure}
 
 The (Fig.\ref{Fig:videostoring}) shows the basic step in which the microscope video is stored. Once the Beaglebone Black has recorded the microscope video, this is upload onto Google App Engine by the board itself, and then the \textit{video flag} is asserted in order to advice the Video Storing application that a new video has been uploaded. Finally, the application downloads and stores this video inside the computer hard drive and resets the flag.
 
 The (List.\ref{code:pro}) shows the taken decision to do not use the graphic user interface (\textit{GUI}) for this application, in order to lighten the application itself, and to allow the use of network.
 
  \begin{figure}[h]
  	\centering
  	\includegraphics[width=\textwidth]{storing/sequence}
  	\caption{Video Storing: Sequence Diagram}
  	\label{Fig:sequenceStoring}
  	
  \end{figure}
\part{Google Glass Application}
\clearpage
\chapter{Wearable Computing - Google Glass }\label{ch:IIIintroduction}
In a world where the electronic technology is growing up faster than every other technology, and where portable devices are strongly common used by people, it is possible to be impressed by a new designed device. Looking to the science fictions, such as \textit{Star Trek}, it was predictable that sooner or later wearable computing is going to be part of everyday life, and so it is happening. Further, this seems to be the next step of the process in how we use computers (Fig.\ref{Fig:grow}), going from desktop usable in a fixed location only, to portability of laptops usable and connected everywhere thank to the wireless connectivity, passing to smartphones and tablets which offer a more portability, to finally wearables computing that, for the time being, offer the highest portability \cite{DDGG}.


\begin{figure}[h]
	\centering
	\includegraphics[width=\textwidth]{GoogleGlass/grow}
	\caption{Computing Evolution}
	\label{Fig:grow}
	
\end{figure}

Thus, \textit{smartglasses} represent an emerging space and not only Google is trying to exploit it but also other famous company like Microsoft and Sony.

\section{Google Glass}

The Google Glass hardware is listed in (Fig.\ref{Fig:glasshw}), and we can find almost all the technology installed in a smartphone, in fact Glass is made with: a battery, a micro-USB port which has the multiple functions of power source, headphone connector, and data port, a bone conduction transducer audio (\textit{BCT}) speaker, a touchpad, different sensors such as the accelerometer and gyroscope to detect head movements, WiFi and Bluetooth interface module, a camera, a microphone, and a prism which acts as display.

Google Glass is not an immersive experience, but it is on only when the user wants it, by watching on the top right corner, where the $640\ x\ 360$ display is placed, and says "\textit{Ok Google}" or tapping onto touchpad.

The Google Glass applications are called \textit{\textbf{Glassware}} and they are available in two different way: Google Mirror API and Glass Development Kit (\textit{GDK}).

\begin{figure}[h]
	\centering
	\includegraphics[width=\textwidth]{GoogleGlass/Hardware.eps}
	\caption{Google Glass}
	\label{Fig:glasshw}
\end{figure}

The \textit{timeline} is a chronological list of cards, each card represents a Glassware, does not matter which kind of it is.
In other words, timeline is a wy to organize the opened application, and users scrolling through different section of the timeline is able to reveal cards in the past, present, and future.

\subsection{Mirror API}

Mirror API Glassware is a web-based application. In fact, the code does not run on the Google Glass itself, but it runs purery in the cloud provided by Google. In this way many programming language are available: Java, Python, PHP, Ruby, .NET and Go.

The main advantage of this applications is that, because the Glassware does not run on Glass, its processor is left free to work on other things and it is not in charge to make some data manipulation, calculation and so on.. 

This service have two big requirements:  of course, being connecting and cloud-aware. Mirror API Glassware has not executable file, and has no access to the Glass sensors.

\subsection{Glass Development Kit}

Glassware built using the Glass Development Kit offers more granular control of application. The programming language available is only one: Java.

This may be defined the classic Android way, where Java code is built in an \textit{APK} (Android Package) file and installed inside the Glass device.

GDK based Glassware has more functionally tha the Mirror API one, indeed it may directly access the hardware, ensure real-time interactions, and being able offline.

The Glassware designed in this thesis is GDK based.

\chapter{The Glassware}

\begin{figure}[h]
	\centering
	\includegraphics[width=\textwidth]{GoogleGlass/interactions}
	\caption{Interactions Between Class}
	\label{Fig:interaction}
\end{figure}

In (Fig.\ref{Fig:interaction}) how the classes and actors in the designed Glassware interact from each others is shown. This is also how a typical \textit{GDK} based Glassware works.



\begin{figure}[b]
	\centering
	\includegraphics[scale =.5]{GoogleGlass/main}
	\caption{MainService's Runnable in Background}
	\label{Fig:mainrun}
\end{figure}

It stars with the service called \textit{MainService} (List.\ref{code:MainService}) that runs when the Glassware voice command is triggered or the application is tapped from the Google Glass Glassware list. When the service has been started, first it creates a live card and then it runs four runnable in background, (Fig.\ref{Fig:mainrun}):
\begin{itemize}
	\item \textit{Sensors Values}, a runnable that periodically downloads and updates the pH and temperature sensors values which then will be plotted.
	\item \textit{Video}, a runnable that periodically checks if a new microscope video is available and, if so, downloads this video inside the Google Glass replacing the older one.
	\item \textit{Beating Image}, a runnable that periodically downloads the beating plot replacing the older one.
	\item \textit{Electrovalves Status}, a runnable that periodically downloads the status of each valve.
\end{itemize}

A \textit{live card} is one of the most important element in the \textit{Google Glass User Interface} (\textit{UI}), it appears in the present section of the timeline and shows information that is relevant at the current time. Live cards can access low level Glass hardware, such as camera, sensors, communication modules, and so on... 

A live card persists in the present section of the timeline as long as the user thinks it may be relevant. Indeed, it is not persistent in the timeline and user can explicitly remove it through a swipe down on the touchpad. Further, more than one live cards may run at the same time, in this case, a live card is still running even if not visible and the focus is on another one.

Live cards are divided in two categories:
\begin{enumerate}
	\item \textit{Low-Frequency Rendering}, limited to a small set of views and can only update the display once every few seconds.
	\item \textit{High-Frequency Rendering}, involve content that is update frequently, more time in a second. This is very useful for plotting data coming from sensors, exactly like the application in which this Glassware is involved, this is why this kind of live card had been used in this thesis project. To create this type of live card an inflating layout has been used.
\end{enumerate}

Finally, live cards exist as long as the Android service which statically generated it is running.

\begin{figure}[h]
	\centering
	\includegraphics[width=\textwidth]{GoogleGlass/Livecard}
	\caption{Live Card, Adding to the Timeline}
	\label{Fig:livecard}
\end{figure}

Backing to the Glassware architecture, once the live card has been created, it constructs an instance of the class \textit{AppDrawer}, an implementation of \textit{Callback} for rendering and thus managing the Glassware \textit{UI}. Live card also adds an instance of the \textit{LiveCardMenu} class, which dynamically manages the menu of the Glassware. Indeed, the application menu changes based on the actual view displayed on the live card.

Then \textit{AppDrawer} creates instances of View subclass, the number of them is four: pH plot, temperature plot, beating plot and electrovalves status. In the (Fig.\ref{Fig:interaction}), these four instance are not shown for compact reasons, instead of them a generic \textit{AppViewer} class is present.

\textit{AppViewer} performs the real rendering, setting the layout and drawing the content on that.

Always in (Fig.\ref{Fig:interaction}) with continuous lines, the steps made when the Glassware has just been opened are shown, while the dashed lines highlight the interactions, that take part when the user acts on Glass.

\section{The Classes}
Below the different classes that compose the Glassware are going to be shown.
\subsection{Main Service}

The code of the main service class is displayed in (List.\ref{code:MainService}). As already said, this class is in charge to create the live card, and to it the live time of the Glassware itself is tied. So, this happens when the service is started at the \textit{OnStartCommand}. After that, an \textit{AppDrawer} instance is created and set as callback in order to render the live card surface.

The \textit{OnStartCommand} method also creates the runnable tasks which run in background to update the data information coming from the Google App Engine Datastore

Finally the main service is also in charge to handle the user requests, that are passed through the menu activity.

\subsection{App Drawer}

The code of the App drawer is shown in (List.\ref{code:AppDrawer}). This subclass allows the control of display surface and chooses which View as to be set at a given time. In fact, it first creates an instance of each view subclass, with their \textit{listener} interface. Then implements the listener method, this is a convenient way to skip from a view to another one. Finally it runs the view draw request, locking the surface canvas.

\subsection{Menu Activity}
The code of the menu activity is shown in (List.\ref{code:MenuActivity}). This subclass plays a very relevant role in the  Glassware functionality. Indeed, it shows the different chooses that allow the user to navigate into the application.

To fulfill that aim it has to change dynamically, based on what is the viewer on focus of the live card. From the main view, user can choose one of the following options: \textit{View pH}, \textit{View Temperature}, \textit{View Video}, \textit{View Beating}, \textit{Drive Electrovalves}, and \textit{Exit}.

\begin{figure}[h]
	\centering
	\includegraphics[width=\textwidth]{GoogleGlass/mainmenu}
	\caption{Main Menu}
	\label{Fig:mainmenu}
	\hfill
\end{figure}

To navigate the menu, user must scroll with a finger, swiping on the touchpad.

The image in (Fig.\ref{Fig:mainmenu}) shows the menu options from the Glassware main view. The (Fig.\ref{Fig:menu1}) shows the menu options when the graphs (pH, temperature, and beating) are displayed.\\
Finally, (Fig.\ref{Fig:menuev}) shows the menu options when the \textit{Drive Electrovalves} is displayed. In that figure only four choses of eight are shown, for compact reasons only. To toggle the valves the Glassware uses two classes: Electrovalves, static class where the status of the valves is stored, and HTTP subclass of Android service, used to send HTTP \textit{GET} and \textit{POST}. The parameter of the HTTP request are attached to the intent that creates the service.



\begin{figure}[h]
	\centering
	\includegraphics[width=0.75\textwidth]{GoogleGlass/menu1}
	\caption{Graph Menu}
	\label{Fig:menu1}
\end{figure}

\begin{figure}[h]
	\centering
	\includegraphics[width=\textwidth]{GoogleGlass/menuev}
	\caption{Menu Electrovalves}
	\label{Fig:menuev}
\end{figure}

Thus, the menu activity is in charge to show dynamically the possible options as well as handle the requests using the Android \textit{intents} to inform the other subclasses what they are supposed to do.

\subsection{App Manager}

The code of this class is shown in (List.\ref{code:AppManager}). The \textit{App Manager} is a static class that is in charge to memorize the state, what is the view that has to be displayed.

\subsection{Main View}

The code of this class is shown in (Lis.\ref{code:MainView}). The \textit{Main View} instantiates the layout \textit{XML} file and sets the Harvard-MIT logo as background image, \textit{"KLab Glassware Interface"} as main written, "Tab to see the menu" as left footer, and the current hour as right footer.
\begin{figure}[h]
	\centering
	\includegraphics[width=.55\textwidth]{GoogleGlass/main}
	\caption{Main View}
	\label{Fig:main}
\end{figure}

\subsection{PH and Temperature Views}

Their codes are shown in (List.\ref{code:PHViewer}) and (List.\ref{code:TemperatureView}). As for the main view, also this subclass instantiates a layout \textit{XML} file and plot the data, write the kind of sensor is it ("\textit{Temperature}" or "\textit{pH}") on the top, and the value of the last data point on the bottom. 
\begin{figure}[h]
	\centering
	\includegraphics[width=.55\textwidth]{GoogleGlass/temperature}
	\caption{Temperature View}
	\label{Fig:temperature}
\end{figure}

\subsection{Beating View}

The last plot, the beating one, is shown in (Fig.\ref{Fig:beating}). The code of this subclass is listed in (List.\ref{code:BeatingView}). In this case the layout is made by an \textit{ImageView} only, and the subclass just set the stored beating plot image.

\begin{figure}[h]
	\centering
	\includegraphics[width=.55\textwidth]{GoogleGlass/beating}
	\caption{Beating View}
	\label{Fig:beating}
\end{figure}

\subsection{Electrovalves View}

\begin{figure}[h]
	\centering
	\subfloat[When Status Has Not Dowloaded Yet\label{subfig-1}]{%
		\includegraphics[width=.5\textwidth]{GoogleGlass/bobo}
	}
	\subfloat[When They Are All Off\label{subfig-2}]{%
		\includegraphics[width=.5\textwidth]{GoogleGlass/electrovalves_all_off}
	}
	\hfill
	\subfloat[When Some of Them is On\label{subfig-3}]{%
		\includegraphics[width=.5\textwidth]{GoogleGlass/elecreovalves_random_status}
	}
	
	
	\caption{\textit{Electrovalves View}}
	\label{elele}
	
\end{figure}

The (Fig.\ref{elele}) shows the result of the electrovalves view. It lists all the eight valves showing the corresponding status. When the user opens the \textit{Drive Electrovalves} tab and the background tasks has not downloaded the valves status yet, the whole electrovalves are displayed without their status, and white written, (Fig.\ref{subfig-1}). When the valves status has been downloaded, the name and the status appear in green color if the valve is on, or in red if it is off, (Fig.\ref{subfig-2}) and (Fig.\ref{subfig-3}).

\subsection{Video Activity}

To play the stored video on Google Glass there are two way, the first one uses the \textit{VideoView} object, and the second one uses an intent belonging to Google Glass API: com.google.glass.action.VIDEOPLAYER. The Glassware described so far exploits this second way, because so the video player is optimized for the Google Glass. Indeed, user can navigate into the video with simple horizontal swipes on the touchpad (Fig.\ref{subfig-45}), can stop the video with a simple tap (Fig.\ref{subfig-35}), and to close the video uses a vertical swipe.

\begin{figure}[h]
	\centering
	\subfloat[Video\label{subfig-45}]{%
		\includegraphics[width=.5\textwidth]{GoogleGlass/video}
	}
	\subfloat[Video on Pause\label{subfig-35}]{%
		\includegraphics[width=.5\textwidth]{GoogleGlass/video_pause}
	}
	
	
	\caption{\textit{Microscope Video}}
	\label{ool}
	
\end{figure}

\clearpage
\part{Conclusion and Appendix}
\chapter{Test and performance}\label{sec:Perf}
\begin{figure}[h]
	\begin{center}
		\includegraphics[width=\textwidth]{breadboard}
		\caption{Circuit mounted on breadboard side view}
		\label{Fig:circuitsfumatp}
	\end{center}
\end{figure}
For what concern the sensors acquiring paths, unfortunately, rigorous experiments have not been done. This is due to the end of my six months period. What I can say for those parts is that from the experiment that I made (not in a microfluidic or biomedical context) the results appear acceptable.

	While, in order to try the reverse control from the Google Glass to the electrovalves we made different kind of experiments.
	
	\section{LED Experiments}
	 First of all we tried the circuit on a breadboard using LEDs instead of electrovalves. The aim of this step is to demonstrate that the firmware running on the Beaglebone Black, the Java code running on the Google Glass and the Python code running on the Google App Engine (used to store the information about the electrovalves status) work well.\\
	Moreover the LED and the electrovalve have basically the same behavior so, if everything works well with the LEDs, there are all the reasons to believe that everything is going to work well with the electrovalves, too.\\
	
	The circuit that actually drives the LEDs is very simple, and it is based on a MOS transistor (\href{http://www.onsemi.com/pub_link/Collateral/BS170-D.PDF}{\textit{BS170}}) used as a switch voltage-controlled, as shown in (Fig.\ref{Fig:driverLED}).
	
	\begin{figure}[h]
		\centering
		\includegraphics[]{Driver/driverLED}
		\caption{Driver for LED}
		\label{Fig:driverLED}
	\end{figure}
	
	The (Fig.\ref{Fig:circuitLED}) shows the circuit used for this step of testing. As can be seen the number of LEDs used is eight, the same number of electrovalves that can be driven from this system.\\
	In order to test all of them we made 2 different kind of trials:
	\begin{enumerate}
		\item \textit{In order turning on\&off}, first all the LEDs are turned on starting from the first one (on the top right corner) to the last one (on the bottom left corner). Then the LEDs are turned off following the same order.\\The result of this can bee watched in  \href{http://youtu.be/iYeAMpxM9uI}{\textbf{this}} video.
		\item \textit{Out of order turning on\&off}, in this trial, like before, all the LEDs start from a condition where all of them are off and then we turned on and off all the LEDs, but in this case following a random order.\\ The result of this can bee watched in  \href{http://youtu.be/mIoylW334Ck}{\textbf{this}} video.
	\end{enumerate}
	
	
	\begin{figure}[h]
		\centering
		\includegraphics[scale=.21]{Experiments/ledBoard}
		\caption{LEDs experiments board}
		\label{Fig:circuitLED}
	\end{figure}
	
	
	\section{Electrovalves experiments}
	
	\subsection{Breadboard Phase}
	The (Fig.\ref{Fig:circuitBreadboard}) shows from the top view the circuit used during the second phase of experiment, the one where we started using electrovalves in a real microfluidic application.\\
	On the left side of the figure we can see the conditioning circuits for the temperature sensor (on the top) and pH sensor (on the bottom). While, on the other side, we can see the part of circuit in charge to drive the electrovalves.\\
	In this last one we are going to focus for now. Each electrovalve is driven by the circuit shown in (Fig.\ref{Fig:driverEV}).
		
	As you can see, this circuit is pretty close to the one of (Fig.\ref{Fig:driverLED}), indeed the only difference is given by freewheeling diode, mandatory because of inductive behavior of electrovalve's solenoid.
	
	\begin{figure}[h]
		\centering
		\includegraphics[scale=.14]{circuitBreadboard}
		\caption{Circuit mounted on breadboard top view}
		\label{Fig:circuitBreadboard}
		
		
	\end{figure}
	
	
	The result of the experiments with electrovalves in a real microfluidic case can bee watched in  \href{https://www.youtube.com/watch?v=CavCVnD2P1k}{\textbf{this}} video.
	
	
	\subsection{PCB Phase}
	
	Finally we replied the last experiment using a PCB (Fig.\ref{fig:PCB}),  designed for this system.\\
	As expected the result of this experiment is the same of the previous step. 
	
	
\chapter{Limitations and Improvements}

As already said, unfortunately, the experiments made to test the sensors acquiring have not be made in a biomedical context.\\
What I did was, in a first time, to test whether the Glassware is able to plot dummy data stored inside the Google App Engine. And in a second time, I tried to plot data from pH and temperature sensors, but not in an organ-on-a-chip application. In both cases the results of experiments were positive, in any case I strongly suggest experiments in a biomedical context.

The embedded system that has been developed involves Internet connection and a server. every time that we are talking about Internet of Think, it is important to deal with Internet security. Thus, the next step of this project regards how to keep the data safe, and how to ensure the privacy, in such a way that no one which is not allowed to see data can have access to them. 

A first, and fast way to fulfill this aim is  using accounts. In fact, \textit{Flask} framework gives the possibility to hide some link if the user is not logged in the web application. In this way, the \textit{Qt} application, the Glassware, and the Embedded Linux firmware have to log-in as the first step.

As is shown in (Fig.\ref{Fig:board}), \textit{PCB} presents two connectors that are not used yet: \textit{Buttons} and \textit{Display} connectors. As can be seen fro the circuit schematic, (Fig.\ref{Fig:circuit}), the first connector has its pins tied to \textit{GPIO} pins of the microcontroller and it is supposed to be used in order to add 4 additional buttons. While the second one is tied to \textit{I2C0} interface of the microcontroller and is though to connect a \textit{I2C display}. In this way, the circuit allows some additional applications that do not require the Google Glass and PC interactions. For example a possible, and really useful, application is to include the sensors calibration. Indeed, for the time being, user has to make the calibration and then modify the bash script which runs the sensor acquiring, (List.\ref{code:sensorbash}):

\begin{lstlisting}[basicstyle=\footnotesize]
./sensor_acquiring [slope_temperature slope_ph 
                    offset_temperature offset_ph]
\end{lstlisting}

While embedding the sensors calibration, user must perform the same calibration steps as before, but the linear regression and transfer function modification  are automatically performed by the system.


\appendix
\chapter{Code}
\definecolor{mygreen}{rgb}{0,0.6,0}
\definecolor{mygray}{rgb}{0.5,0.5,0.5}
\definecolor{mymauve}{rgb}{0.58,0,0.82}

\section{Firmware}\label{code:firmware}

\lstinputlisting[frame=single,  basicstyle=\footnotesize, breakatwhitespace=true, 
basicstyle=\scriptsize\ttfamily,,
breaklines=false,
commentstyle=\color{mygreen},
keepspaces=true,
keywordstyle=\color{blue}, 
numbers=left,
numbersep=5pt, 
numberstyle=\tiny\color{mygray}, % the style that is used for the line-numbers
rulecolor=\color{black}, 
showspaces=false,
showstringspaces=false,          % underline spaces within strings only
showtabs=false,                  % show tabs within strings adding particular underscores
stepnumber=2,                    % the step between two line-numbers. If it's 1, each line will be numbered
stringstyle=\color{mymauve},     % string literal style
tabsize=2,                       % sets default tabsize to 2 spaces
language=bash,
caption = KlabFirware (bash),
label = code:recordVideo]{./code/beaglebone/Video/recordVideo}

\clearpage

\lstinputlisting[frame=single,  basicstyle=\footnotesize, breakatwhitespace=true, 
basicstyle=\scriptsize\ttfamily,,
breaklines=false,
commentstyle=\color{mygreen},
keepspaces=true,
keywordstyle=\color{blue}, 
numbers=left,
numbersep=5pt, 
numberstyle=\tiny\color{mygray}, % the style that is used for the line-numbers
rulecolor=\color{black}, 
showspaces=false,
showstringspaces=false,          % underline spaces within strings only
showtabs=false,                  % show tabs within strings adding particular underscores
stepnumber=2,                    % the step between two line-numbers. If it's 1, each line will be numbered
stringstyle=\color{mymauve},     % string literal style
tabsize=2,                       % sets default tabsize to 2 spaces
language=bash,
caption = SensorAcquiring (bash),
label = code:sensorbash]{./code/beaglebone/Sensors/sensorAcquiring}

\lstinputlisting[frame=single,  basicstyle=\footnotesize, breakatwhitespace=true, 
basicstyle=\scriptsize\ttfamily,,
breaklines=false,
commentstyle=\color{mygreen},
keepspaces=true,
keywordstyle=\color{blue}, 
numbers=left,
numbersep=5pt, 
numberstyle=\tiny\color{mygray}, % the style that is used for the line-numbers
rulecolor=\color{black}, 
showspaces=false,
showstringspaces=false,          % underline spaces within strings only
showtabs=false,                  % show tabs within strings adding particular underscores
stepnumber=2,                    % the step between two line-numbers. If it's 1, each line will be numbered
stringstyle=\color{mymauve},     % string literal style
tabsize=2,                       % sets default tabsize to 2 spaces
language=C++,
caption = Pins Setting,
label = code:pins]{./code/beaglebone/Pins/bo.dts}

\lstinputlisting[frame=single,  basicstyle=\footnotesize, breakatwhitespace=true, 
basicstyle=\scriptsize\ttfamily,,
breaklines=false,
commentstyle=\color{mygreen},
keepspaces=true,
keywordstyle=\color{blue}, 
numbers=left,
numbersep=5pt, 
numberstyle=\tiny\color{mygray}, % the style that is used for the line-numbers
rulecolor=\color{black}, 
showspaces=false,
showstringspaces=false,          % underline spaces within strings only
showtabs=false,                  % show tabs within strings adding particular underscores
stepnumber=2,                    % the step between two line-numbers. If it's 1, each line will be numbered
stringstyle=\color{mymauve},     % string literal style
tabsize=2,                       % sets default tabsize to 2 spaces
language=C,
caption = ElectrovalvesDriver.cpp,
label = code:Electrovalves]{./code/beaglebone/Electrovalves/GlassInterface.cpp}

\lstinputlisting[frame=single,  basicstyle=\footnotesize, breakatwhitespace=true, 
basicstyle=\scriptsize\ttfamily,,
breaklines=false,
commentstyle=\color{mygreen},
keepspaces=true,
keywordstyle=\color{blue}, 
numbers=left,
numbersep=5pt, 
numberstyle=\tiny\color{mygray}, % the style that is used for the line-numbers
rulecolor=\color{black}, 
showspaces=false,
showstringspaces=false,          % underline spaces within strings only
showtabs=false,                  % show tabs within strings adding particular underscores
stepnumber=2,                    % the step between two line-numbers. If it's 1, each line will be numbered
stringstyle=\color{mymauve},     % string literal style
tabsize=2,                       % sets default tabsize to 2 spaces
language=C++,
caption = sensor\_aquiring.cpp,
label = code:sensors]{./code/beaglebone/Sensors/sensor_acquiring.cpp}



\lstinputlisting[frame=single,  basicstyle=\footnotesize, breakatwhitespace=true, 
basicstyle=\scriptsize\ttfamily,,
breaklines=false,
commentstyle=\color{mygreen},
keepspaces=true,
keywordstyle=\color{blue}, 
numbers=left,
numbersep=5pt, 
numberstyle=\tiny\color{mygray}, % the style that is used for the line-numbers
rulecolor=\color{black}, 
showspaces=false,
showstringspaces=false,          % underline spaces within strings only
showtabs=false,                  % show tabs within strings adding particular underscores
stepnumber=2,                    % the step between two line-numbers. If it's 1, each line will be numbered
stringstyle=\color{mymauve},     % string literal style
tabsize=2,                       % sets default tabsize to 2 spaces
language=C++,
caption = compute\_beating.cpp,
label = code:beating]{./code/beaglebone/Video/compute_beating.cpp}

\lstinputlisting[frame=single,  basicstyle=\footnotesize, breakatwhitespace=true, 
basicstyle=\scriptsize\ttfamily,,
breaklines=false,
commentstyle=\color{mygreen},
keepspaces=true,
keywordstyle=\color{blue}, 
numbers=left,
numbersep=5pt, 
numberstyle=\tiny\color{mygray}, % the style that is used for the line-numbers
rulecolor=\color{black}, 
showspaces=false,
showstringspaces=false,          % underline spaces within strings only
showtabs=false,                  % show tabs within strings adding particular underscores
stepnumber=2,                    % the step between two line-numbers. If it's 1, each line will be numbered
stringstyle=\color{mymauve},     % string literal style
tabsize=2,                       % sets default tabsize to 2 spaces
language=C++,
caption = Pins Setting,
label = code:pins]{./code/beaglebone/Pins/bo.dts}

%\lstinputlisting[frame=single,  basicstyle=\footnotesize, breakatwhitespace=true, 
%basicstyle=\scriptsize\ttfamily,,
%breaklines=false,
%commentstyle=\color{mygreen},
%keepspaces=true,
%keywordstyle=\color{blue}, 
%numbers=left,
%numbersep=5pt, 
%numberstyle=\tiny\color{mygray}, % the style that is used for the line-numbers
%rulecolor=\color{black}, 
%showspaces=false,
%showstringspaces=false,          % underline spaces within strings only
%showtabs=false,                  % show tabs within strings adding particular underscores
%stepnumber=2,                    % the step between two line-numbers. If it's 1, each line will be numbered
%stringstyle=\color{mymauve},     % string literal style
%tabsize=2,                       % sets default tabsize to 2 spaces
%language=C++,
%caption = GPIO.h,
%label = code:gpio]{./code/beaglebone/GPIO.h}
%
%\lstinputlisting[frame=single,  basicstyle=\footnotesize, breakatwhitespace=true, 
%basicstyle=\scriptsize\ttfamily,,
%breaklines=false,
%commentstyle=\color{mygreen},
%keepspaces=true,
%keywordstyle=\color{blue}, 
%numbers=left,
%numbersep=5pt, 
%numberstyle=\tiny\color{mygray}, % the style that is used for the line-numbers
%rulecolor=\color{black}, 
%showspaces=false,
%showstringspaces=false,          % underline spaces within strings only
%showtabs=false,                  % show tabs within strings adding particular underscores
%stepnumber=2,                    % the step between two line-numbers. If it's 1, each line will be numbered
%stringstyle=\color{mymauve},     % string literal style
%tabsize=2,                       % sets default tabsize to 2 spaces
%language=C++,
%caption = GPIO.cpp]{./code/beaglebone/GPIO.cpp}

\lstinputlisting[frame=single,  basicstyle=\footnotesize, breakatwhitespace=true, 
basicstyle=\scriptsize\ttfamily,,
breaklines=false,
commentstyle=\color{mygreen},
keepspaces=true,
keywordstyle=\color{blue}, 
numbers=left,
numbersep=5pt, 
numberstyle=\tiny\color{mygray}, % the style that is used for the line-numbers
rulecolor=\color{black}, 
showspaces=false,
showstringspaces=false,          % underline spaces within strings only
showtabs=false,                  % show tabs within strings adding particular underscores
stepnumber=2,                    % the step between two line-numbers. If it's 1, each line will be numbered
stringstyle=\color{mymauve},     % string literal style
tabsize=2,                       % sets default tabsize to 2 spaces
language=C++,
caption = HTTP.h,
label = code:http]{./code/beaglebone/http.h}

\lstinputlisting[frame=single,  basicstyle=\footnotesize, breakatwhitespace=true, 
basicstyle=\scriptsize\ttfamily,,
breaklines=false,
commentstyle=\color{mygreen},
keepspaces=true,
keywordstyle=\color{blue}, 
numbers=left,
numbersep=5pt, 
numberstyle=\tiny\color{mygray}, % the style that is used for the line-numbers
rulecolor=\color{black}, 
showspaces=false,
showstringspaces=false,          % underline spaces within strings only
showtabs=false,                  % show tabs within strings adding particular underscores
stepnumber=2,                    % the step between two line-numbers. If it's 1, each line will be numbered
stringstyle=\color{mymauve},     % string literal style
tabsize=2,                       % sets default tabsize to 2 spaces
language=C++,
caption = HTTP.cpp]{./code/beaglebone/http.cpp}

\section{Video Storing Software}\label{code:video}

\lstinputlisting[frame=single,  basicstyle=\footnotesize, breakatwhitespace=true, 
basicstyle=\scriptsize\ttfamily,,
breaklines=false,
commentstyle=\color{mygreen},
keepspaces=true,
keywordstyle=\color{blue}, 
numbers=left,
numbersep=5pt, 
numberstyle=\tiny\color{mygray}, % the style that is used for the line-numbers
rulecolor=\color{black}, 
showspaces=false,
showstringspaces=false,          % underline spaces within strings only
showtabs=false,                  % show tabs within strings adding particular underscores
stepnumber=2,                    % the step between two line-numbers. If it's 1, each line will be numbered
stringstyle=\color{mymauve},     % string literal style
tabsize=2,                       % sets default tabsize to 2 spaces
language=C++,
caption = main.cpp,
label = code:mainstoring]{./code/storing/main.cpp}

\lstinputlisting[frame=single,  basicstyle=\footnotesize, breakatwhitespace=true, 
basicstyle=\scriptsize\ttfamily,,
breaklines=false,
commentstyle=\color{mygreen},
keepspaces=true,
keywordstyle=\color{blue}, 
numbers=left,
numbersep=5pt, 
numberstyle=\tiny\color{mygray}, % the style that is used for the line-numbers
rulecolor=\color{black}, 
showspaces=false,
showstringspaces=false,          % underline spaces within strings only
showtabs=false,                  % show tabs within strings adding particular underscores
stepnumber=2,                    % the step between two line-numbers. If it's 1, each line will be numbered
stringstyle=\color{mymauve},     % string literal style
tabsize=2,                       % sets default tabsize to 2 spaces
language=C++,
caption = VideoStoring.pro,
label = code:pro]{./code/storing/VideoStoring.pro}

\lstinputlisting[frame=single,  basicstyle=\footnotesize, breakatwhitespace=true, 
basicstyle=\scriptsize\ttfamily,,
breaklines=false,
commentstyle=\color{mygreen},
keepspaces=true,
keywordstyle=\color{blue}, 
numbers=left,
numbersep=5pt, 
numberstyle=\tiny\color{mygray}, % the style that is used for the line-numbers
rulecolor=\color{black}, 
showspaces=false,
showstringspaces=false,          % underline spaces within strings only
showtabs=false,                  % show tabs within strings adding particular underscores
stepnumber=2,                    % the step between two line-numbers. If it's 1, each line will be numbered
stringstyle=\color{mymauve},     % string literal style
tabsize=2,                       % sets default tabsize to 2 spaces
language=C++,
caption = mytimer.h,
label = code:timer]{./code/storing/mytimer.h}

\lstinputlisting[frame=single,  basicstyle=\footnotesize, breakatwhitespace=true, 
basicstyle=\scriptsize\ttfamily,,
breaklines=false,
commentstyle=\color{mygreen},
keepspaces=true,
keywordstyle=\color{blue}, 
numbers=left,
numbersep=5pt, 
numberstyle=\tiny\color{mygray}, % the style that is used for the line-numbers
rulecolor=\color{black}, 
showspaces=false,
showstringspaces=false,          % underline spaces within strings only
showtabs=false,                  % show tabs within strings adding particular underscores
stepnumber=2,                    % the step between two line-numbers. If it's 1, each line will be numbered
stringstyle=\color{mymauve},     % string literal style
tabsize=2,                       % sets default tabsize to 2 spaces
language=C++,
caption = mytimer.cpp]{./code/storing/mytimer.cpp}

\lstinputlisting[frame=single,  basicstyle=\footnotesize, breakatwhitespace=true, 
basicstyle=\scriptsize\ttfamily,,
breaklines=false,
commentstyle=\color{mygreen},
keepspaces=true,
keywordstyle=\color{blue}, 
numbers=left,
numbersep=5pt, 
numberstyle=\tiny\color{mygray}, % the style that is used for the line-numbers
rulecolor=\color{black}, 
showspaces=false,
showstringspaces=false,          % underline spaces within strings only
showtabs=false,                  % show tabs within strings adding particular underscores
stepnumber=2,                    % the step between two line-numbers. If it's 1, each line will be numbered
stringstyle=\color{mymauve},     % string literal style
tabsize=2,                       % sets default tabsize to 2 spaces
language=C++,
caption = downloader.h,
label = code:downloader]{./code/storing/downloader.h}

\lstinputlisting[frame=single,  basicstyle=\footnotesize, breakatwhitespace=true, 
basicstyle=\scriptsize\ttfamily,,
breaklines=false,
commentstyle=\color{mygreen},
keepspaces=true,
keywordstyle=\color{blue}, 
numbers=left,
numbersep=5pt, 
numberstyle=\tiny\color{mygray}, % the style that is used for the line-numbers
rulecolor=\color{black}, 
showspaces=false,
showstringspaces=false,          % underline spaces within strings only
showtabs=false,                  % show tabs within strings adding particular underscores
stepnumber=2,                    % the step between two line-numbers. If it's 1, each line will be numbered
stringstyle=\color{mymauve},     % string literal style
tabsize=2,                       % sets default tabsize to 2 spaces
language=C++,
caption = downloader.cpp]{./code/storing/downloader.cpp}

\section{Google App Engine}\label{server}

\lstinputlisting[frame=single,  basicstyle=\footnotesize, breakatwhitespace=true, 
basicstyle=\scriptsize\ttfamily,,
breaklines=false,
commentstyle=\color{mygreen},
keepspaces=true,
keywordstyle=\color{blue}, 
numbers=left,
numbersep=5pt, 
numberstyle=\tiny\color{mygray}, % the style that is used for the line-numbers
rulecolor=\color{black}, 
showspaces=false,
showstringspaces=false,          % underline spaces within strings only
showtabs=false,                  % show tabs within strings adding particular underscores
stepnumber=2,                    % the step between two line-numbers. If it's 1, each line will be numbered
stringstyle=\color{mymauve},     % string literal style
tabsize=2,                       % sets default tabsize to 2 spaces
language=python,
caption = Main Script of Server,
label = code:server]{./code/main.py}

\section{Glassware}

\lstinputlisting[frame=single,  basicstyle=\footnotesize, breakatwhitespace=true, 
basicstyle=\scriptsize\ttfamily,,
breaklines=false,
commentstyle=\color{mygreen},
keepspaces=true,
keywordstyle=\color{blue}, 
numbers=left,
numbersep=5pt, 
numberstyle=\tiny\color{mygray}, % the style that is used for the line-numbers
rulecolor=\color{black}, 
showspaces=false,
showstringspaces=false,          % underline spaces within strings only
showtabs=false,                  % show tabs within strings adding particular underscores
stepnumber=2,                    % the step between two line-numbers. If it's 1, each line will be numbered
stringstyle=\color{mymauve},     % string literal style
tabsize=2,                       % sets default tabsize to 2 spaces
language=Java,
caption = MainService.java,
label = code:MainService]{./code/glass/MainService.java}


\lstinputlisting[frame=single,  basicstyle=\footnotesize, breakatwhitespace=true, 
basicstyle=\scriptsize\ttfamily,,
breaklines=false,
commentstyle=\color{mygreen},
keepspaces=true,
keywordstyle=\color{blue}, 
numbers=left,
numbersep=5pt, 
numberstyle=\tiny\color{mygray}, % the style that is used for the line-numbers
rulecolor=\color{black}, 
showspaces=false,
showstringspaces=false,          % underline spaces within strings only
showtabs=false,                  % show tabs within strings adding particular underscores
stepnumber=2,                    % the step between two line-numbers. If it's 1, each line will be numbered
stringstyle=\color{mymauve},     % string literal style
tabsize=2,                       % sets default tabsize to 2 spaces
language=Java,
caption = AppDrawer.java,
label = code:AppDrawer]{./code/glass/AppDrawer.java}

\lstinputlisting[frame=single,  basicstyle=\footnotesize, breakatwhitespace=true, 
basicstyle=\scriptsize\ttfamily,,
breaklines=false,
commentstyle=\color{mygreen},
keepspaces=true,
keywordstyle=\color{blue}, 
numbers=left,
numbersep=5pt, 
numberstyle=\tiny\color{mygray}, % the style that is used for the line-numbers
rulecolor=\color{black}, 
showspaces=false,
showstringspaces=false,          % underline spaces within strings only
showtabs=false,                  % show tabs within strings adding particular underscores
stepnumber=2,                    % the step between two line-numbers. If it's 1, each line will be numbered
stringstyle=\color{mymauve},     % string literal style
tabsize=2,                       % sets default tabsize to 2 spaces
language=Java,
caption = MainView.java,
label = code:MainView]{./code/glass/MainView.java}

\lstinputlisting[frame=single,  basicstyle=\footnotesize, breakatwhitespace=true, 
basicstyle=\scriptsize\ttfamily,,
breaklines=false,
commentstyle=\color{mygreen},
keepspaces=true,
keywordstyle=\color{blue}, 
numbers=left,
numbersep=5pt, 
numberstyle=\tiny\color{mygray}, % the style that is used for the line-numbers
rulecolor=\color{black}, 
showspaces=false,
showstringspaces=false,          % underline spaces within strings only
showtabs=false,                  % show tabs within strings adding particular underscores
stepnumber=2,                    % the step between two line-numbers. If it's 1, each line will be numbered
stringstyle=\color{mymauve},     % string literal style
tabsize=2,                       % sets default tabsize to 2 spaces
language=Java,
caption = BeatingView.java,
label = code:BeatingView]{./code/glass/BeatingView.java}

\lstinputlisting[frame=single,  basicstyle=\footnotesize, breakatwhitespace=true, 
basicstyle=\scriptsize\ttfamily,,
breaklines=false,
commentstyle=\color{mygreen},
keepspaces=true,
keywordstyle=\color{blue}, 
numbers=left,
numbersep=5pt, 
numberstyle=\tiny\color{mygray}, % the style that is used for the line-numbers
rulecolor=\color{black}, 
showspaces=false,
showstringspaces=false,          % underline spaces within strings only
showtabs=false,                  % show tabs within strings adding particular underscores
stepnumber=2,                    % the step between two line-numbers. If it's 1, each line will be numbered
stringstyle=\color{mymauve},     % string literal style
tabsize=2,                       % sets default tabsize to 2 spaces
language=Java,
caption = PHViewer.java,
label = code:PHViewer]{./code/glass/PHViewer.java}

\lstinputlisting[frame=single,  basicstyle=\footnotesize, breakatwhitespace=true, 
basicstyle=\scriptsize\ttfamily,,
breaklines=false,
commentstyle=\color{mygreen},
keepspaces=true,
keywordstyle=\color{blue}, 
numbers=left,
numbersep=5pt, 
numberstyle=\tiny\color{mygray}, % the style that is used for the line-numbers
rulecolor=\color{black}, 
showspaces=false,
showstringspaces=false,          % underline spaces within strings only
showtabs=false,                  % show tabs within strings adding particular underscores
stepnumber=2,                    % the step between two line-numbers. If it's 1, each line will be numbered
stringstyle=\color{mymauve},     % string literal style
tabsize=2,                       % sets default tabsize to 2 spaces
language=Java,
caption = TemperatureView.java,
label = code:TemperatureView]{./code/glass/TemperatureView.java}

\lstinputlisting[frame=single,  basicstyle=\footnotesize, breakatwhitespace=true, 
basicstyle=\scriptsize\ttfamily,,
breaklines=false,
commentstyle=\color{mygreen},
keepspaces=true,
keywordstyle=\color{blue}, 
numbers=left,
numbersep=5pt, 
numberstyle=\tiny\color{mygray}, % the style that is used for the line-numbers
rulecolor=\color{black}, 
showspaces=false,
showstringspaces=false,          % underline spaces within strings only
showtabs=false,                  % show tabs within strings adding particular underscores
stepnumber=2,                    % the step between two line-numbers. If it's 1, each line will be numbered
stringstyle=\color{mymauve},     % string literal style
tabsize=2,                       % sets default tabsize to 2 spaces
language=Java,
caption = VideoPlayerActivity.java,
label = code:VideoPlayerActivity]{./code/glass/VideoPlayerActivity.java}

\lstinputlisting[frame=single,  basicstyle=\footnotesize, breakatwhitespace=true, 
basicstyle=\scriptsize\ttfamily,,
breaklines=false,
commentstyle=\color{mygreen},
keepspaces=true,
keywordstyle=\color{blue}, 
numbers=left,
numbersep=5pt, 
numberstyle=\tiny\color{mygray}, % the style that is used for the line-numbers
rulecolor=\color{black}, 
showspaces=false,
showstringspaces=false,          % underline spaces within strings only
showtabs=false,                  % show tabs within strings adding particular underscores
stepnumber=2,                    % the step between two line-numbers. If it's 1, each line will be numbered
stringstyle=\color{mymauve},     % string literal style
tabsize=2,                       % sets default tabsize to 2 spaces
language=Java,
caption = MenuActivity.java,
label = code:MenuActivity]{./code/glass/MenuActivity.java}

\lstinputlisting[frame=single,  basicstyle=\footnotesize, breakatwhitespace=true, 
basicstyle=\scriptsize\ttfamily,,
breaklines=false,
commentstyle=\color{mygreen},
keepspaces=true,
keywordstyle=\color{blue}, 
numbers=left,
numbersep=5pt, 
numberstyle=\tiny\color{mygray}, % the style that is used for the line-numbers
rulecolor=\color{black}, 
showspaces=false,
showstringspaces=false,          % underline spaces within strings only
showtabs=false,                  % show tabs within strings adding particular underscores
stepnumber=2,                    % the step between two line-numbers. If it's 1, each line will be numbered
stringstyle=\color{mymauve},     % string literal style
tabsize=2,                       % sets default tabsize to 2 spaces
language=Java,
caption = AppManager.java,
label = code:AppManager]{./code/glass/AppManager.java}

\lstinputlisting[frame=single,  basicstyle=\footnotesize, breakatwhitespace=true, 
basicstyle=\scriptsize\ttfamily,,
breaklines=false,
commentstyle=\color{mygreen},
keepspaces=true,
keywordstyle=\color{blue}, 
numbers=left,
numbersep=5pt, 
numberstyle=\tiny\color{mygray}, % the style that is used for the line-numbers
rulecolor=\color{black}, 
showspaces=false,
showstringspaces=false,          % underline spaces within strings only
showtabs=false,                  % show tabs within strings adding particular underscores
stepnumber=2,                    % the step between two line-numbers. If it's 1, each line will be numbered
stringstyle=\color{mymauve},     % string literal style
tabsize=2,                       % sets default tabsize to 2 spaces
language=Java,
caption = DataPoint.java,
label = code:DataPoint]{./code/glass/DataPoint.java}





\cite{SNOA529A}

\clearpage
\listoffigures
\addcontentsline{toc}{chapter}{List of Figures}
\clearpage
\chapter*{list of listings}
\addcontentsline{toc}{chapter}{List of Listings}
\lolnoheading
%\appendix

%\section{Firmware}\label{code:firmware}

\section{Video Storing Software}\label{code:video}

\section{Google App Engine}\label{server}

\section{Glassware}





%\vfill
%\chapter*{BIBLIOGRAPHY}
%\addcontentsline{toc}{chapter}{BIBLIOGRAPHY}
%\markboth{Fabio Busignani Master Thesis}{BIBLIOGRAPHY}
\bibliographystyle{unsrtnat}
\bibliography{biblio}
\addcontentsline{toc}{chapter}{Bibliography}

%BibTex!!!

\end{document}
